%% ============================================================
%%  DIEM Chatbot — Relazione Tecnica
%%  Corso di Natural Language Processing
%%  Università degli Studi di Salerno
%% ============================================================
\documentclass[11pt,a4paper]{article}

\usepackage[utf8]{inputenc}
\usepackage[T1]{fontenc}
\usepackage[italian]{babel}
\usepackage[margin=1.9cm]{geometry}
\usepackage{graphicx}
\usepackage[colorlinks=true,linkcolor=blue!70!black,urlcolor=blue!70!black,citecolor=green!50!black]{hyperref}
\usepackage[table]{xcolor}
\usepackage{colortbl}
\usepackage{booktabs}
\usepackage{tabularx}
\usepackage{longtable}
\usepackage{float}
\usepackage{enumitem}
\usepackage{amsmath,amssymb}
\usepackage[most]{tcolorbox}
\usepackage{titlesec}
\usepackage{fancyhdr}
\usepackage{microtype}
\usepackage{parskip}
\setlength{\parskip}{4pt plus 1pt minus 1pt}
\usepackage{listings}
\usepackage{tikz}
\usetikzlibrary{shapes.geometric,arrows.meta,positioning,fit,backgrounds}
\usepackage{pgfplots}
\pgfplotsset{compat=1.18}
\usepackage{caption}
\usepackage{subcaption}
\usepackage{natbib}
\usepackage{url}
\usepackage{lmodern}

%% --- List compaction ---
\setlist[itemize]{nosep,leftmargin=*,topsep=3pt,partopsep=0pt}
\setlist[enumerate]{nosep,leftmargin=*,topsep=3pt,partopsep=0pt}

%% --- Widows & orphans ---
\widowpenalty=10000
\clubpenalty=10000

%% --- Line breaking tolerance (prevents overfull with long \texttt/URLs) ---
\tolerance=1500
\emergencystretch=2em

%% --- Section styling ---
\titleformat{\section}{\large\bfseries\color{blue!70!black}}{\thesection}{0.8em}{}
\titleformat{\subsection}{\normalsize\bfseries\color{blue!50!black}}{\thesubsection}{0.8em}{}
\titleformat{\subsubsection}{\normalsize\itshape}{\thesubsubsection}{0.8em}{}
\titlespacing*{\section}{0pt}{12pt}{4pt}
\titlespacing*{\subsection}{0pt}{8pt}{3pt}
\titlespacing*{\subsubsection}{0pt}{6pt}{2pt}

%% --- Headers / Footers ---
\pagestyle{fancy}
\fancyhf{}
\fancyhead[L]{\small\textit{DIEM Chatbot --- Sistema RAG Agentico}}
\fancyhead[R]{\small\textit{NLP --- UniSA 2025/2026}}
\fancyfoot[C]{\thepage}
\renewcommand{\headrulewidth}{0.4pt}

%% --- tcolorbox styles ---
\tcbset{
  boxstyle/.style={
    colback=blue!3,colframe=blue!60!black,
    fonttitle=\bfseries\small,
    boxrule=0.5pt,arc=3pt,
    left=6pt,right=6pt,top=4pt,bottom=4pt
  },
  warnstyle/.style={
    colback=orange!8,colframe=orange!70!black,
    fonttitle=\bfseries\small,
    boxrule=0.5pt,arc=3pt
  },
  greenbox/.style={
    colback=green!5,colframe=green!60!black,
    fonttitle=\bfseries\small,
    boxrule=0.5pt,arc=3pt
  },
  s2estyle/.style={
    colback=gray!5,colframe=gray!45!black,
    fonttitle=\bfseries\small\color{gray!60!black},
    boxrule=0.4pt,arc=3pt,
    left=5pt,right=5pt,top=3pt,bottom=3pt,
    breakable
  }
}

%% ---- Note di sviluppo (start-to-end helper) ----
%% Imposta \showsdevnotesfalse prima della consegna finale per nascondere tutti i box S2E.
\newif\ifshowsdevnotes
\showsdevnotestrue
\newcommand{\snote}[1]{%
  \ifshowsdevnotes
  \begin{tcolorbox}[s2estyle,title={Note di sviluppo (start\,\texorpdfstring{$\to$}{->}\,end)}]
  \footnotesize #1
  \end{tcolorbox}
  \fi
}

%% --- Code listing style ---
\lstset{
  basicstyle=\ttfamily\scriptsize,
  keywordstyle=\bfseries\color{blue!70!black},
  commentstyle=\itshape\color{gray},
  stringstyle=\color{green!50!black},
  breaklines=true,
  frame=single,
  frameround=tttt,
  framesep=3pt,
  xleftmargin=6pt,xrightmargin=6pt,
  captionpos=b,
  aboveskip=4pt,
  belowskip=2pt
}

%% --- Compact table rows / itemize spacing for page budget ---
\renewcommand{\arraystretch}{0.90}
\setlength{\abovecaptionskip}{4pt}
\setlength{\belowcaptionskip}{2pt}

%% --- Misc macros ---
\newcommand{\code}[1]{\texttt{\small#1}}
\newcommand{\sys}{\textsc{DiemBot}}

%% ============================================================
\begin{document}

%% ---- TITLE PAGE ----
\begin{titlepage}
  \centering
  \vspace*{2cm}
  {\Large Università degli Studi di Salerno\\[4pt]
   \normalsize Dipartimento di Ingegneria dell'Informazione ed Elettrica e Matematica Applicata}\\[1cm]
  \rule{\textwidth}{1pt}\\[0.5cm]
  {\LARGE\bfseries\color{blue!70!black}
    DIEM Chatbot: Progettazione e Sviluppo di un\\[6pt]
    Sistema RAG Agentico per il Supporto\\[6pt]
    Informativo Universitario
  }\\[0.4cm]
  \rule{\textwidth}{1pt}\\[1.5cm]
  {\large Relazione Tecnica --- Corso di Natural Language Processing}\\[2cm]
  \begin{tabular}{lll}
    \textbf{Autori:} & Alessio Bosco    & \small a.bosco22@studenti.unisa.it \ -- \ 0622702471 \\[3pt]
                     & Michele Apuzzo   & \small m.apuzzo13@studenti.unisa.it \ -- \ 0622702461 \\[3pt]
                     & Andrey Kiriyanov & \small a.kiriyanov@studenti.unisa.it \ -- \ 0622702577 \\[3pt]
                      & Angelo Annunziata & \small a.annunziata111@studenti.unisa.it \ -- \ 0622702489 \\\\[6pt]
    \textbf{Anno Accademico:} & 2025/2026 & \\[4pt]
  \end{tabular}\\[3cm]
\end{titlepage}

\tableofcontents

%% ============================================================
\newpage
\section{Introduzione}
\label{sec:intro}

Gli atenei italiani erogano un volume crescente di informazioni attraverso portali web
frammentati, spesso difficili da navigare per studenti e personale: programmi di corso,
normative di ammissione, schede docente, calendario accademico e informazioni sui
laboratori sono distribuite su domini separati, con strutture disomogenee e contenuti
non sempre aggiornati in modo coerente. La ricerca di un'informazione specifica richiede
spesso più passaggi manuali tra pagine diverse, con rischio di dati obsoleti o
incompleti.

\sys{} nasce per rispondere a questa esigenza nel contesto specifico del Dipartimento di
Ingegneria dell'Informazione ed Elettrica e Matematica Applicata (DIEM) dell'Università
degli Studi di Salerno. L'obiettivo è un assistente conversazionale capace di rispondere
in italiano (e in inglese) a domande fattuali sul dipartimento --- corsi di laurea, piani di
studio, docenti, strutture, procedure amministrative --- sfruttando come knowledge base
i contenuti reali del sito istituzionale.

\subsection{Motivazioni e Sfide}

Il dominio applicativo presenta alcune sfide concrete che hanno guidato le scelte
progettuali. \textbf{Eterogeneità delle fonti:} le informazioni sono distribuite su tre domini
(\texttt{diem.unisa.it}, \texttt{docenti.unisa.it}, \texttt{corsi.unisa.it}) con strutture di navigazione diverse rendendo necessarie strategie di crawling
differenziate. \textbf{Assenza di metadati temporali:} i CMS istituzionali non espongono
segnali di data leggibili automaticamente, il che complica qualsiasi forma
di recency scoring basata su documenti. \textbf{Ambiguità semantica dei chunk:} frammenti
testuali come ``modalità d'esame'', ``ricevimento'' o ``obiettivi formativi'' sono identici
tra decine di pagine diverse: senza informazioni di contesto negli embedding, il retrieval
produce falsi positivi sistematici.

A queste sfide si aggiungono requisiti di sistema tipici di un assistente universitario
reale: rifiuto di domande fuori ambito, gestione di conversazioni multi-turno con
risoluzione di pronomi, redazione di dati sensibili (indirizzi e-mail, numeri di carta),
robustezza ad adversarial input e capacità di calcolo accademico (voti, medie TOLC).

\subsection{Contributi}

Il lavoro presenta i seguenti contributi originali.

\begin{enumerate}
  \item Una pipeline di ingestion end-to-end specifica per siti universitari italiani, con
        crawling deterministico via sitemap, filtraggio temporale e linguistico, arricchimento
        contestuale via LLM (context header) e chunking Parent-Child con intestazioni
        semantiche incorporate in ogni chunk figlio.

  \item Un'architettura RAG agentica basata su LangGraph StateGraph con otto nodi
        specializzati, guardrail a cascata, query rewriting automatico con iniezione
        dell'anno accademico e ciclo di retrieval iterativo controllato da un contatore di
        stato.

  \item Una validazione sperimentale intermedia che ha guidato la stabilizzazione del
        sistema, trasformando errori osservati nelle prime run (retrieval saltato,
        drift multi-turno, rewrite generico, knowledge gap) in interventi mirati su grafo,
        prompt e valutazione.

  \item Un framework di valutazione su golden set italiano con sei metriche Ragas e due
        criteri custom (scope-awareness e robustness).
\end{enumerate}

\subsection{Struttura del documento}

La Sezione~\ref{sec:background} introduce i concetti fondamentali e il lavoro correlato.
La Sezione~\ref{sec:architettura} descrive l'architettura generale del sistema. Le
Sezioni~\ref{sec:ingestion} e~\ref{sec:core-rag} dettagliano rispettivamente la pipeline
di ingestion e il core RAG agentico, inclusi tool, prompt engineering e raffinamenti
del ciclo rewrite--retrieve. La Sezione~\ref{sec:evaluation} presenta il framework di
valutazione, mentre la Sezione~\ref{sec:stabilization} descrive la validazione intermedia
e le decisioni di stabilizzazione. La Sezione~\ref{sec:eval_results} analizza i risultati
quantitativi sulle configurazioni C1--C5. Le Sezioni~\ref{sec:discussion}
e~\ref{sec:conclusioni} discutono i limiti, i lavori futuri e le conclusioni del lavoro.

%% ============================================================
\section{Background e Lavori Correlati}
\label{sec:background}

Il paradigma Retrieval-Augmented Generation (RAG) \citep{lewis2020rag} supera un limite
fondamentale degli LLM puri: la conoscenza cristallizzata nel training non può essere
aggiornata senza un costoso re-fine-tuning. Il modello generativo riceve, oltre alla query,
un contesto recuperato dinamicamente da una knowledge base esterna, riducendo le
allucinazioni e permettendo aggiornamento incrementale e citazione delle fonti.
L'architettura classica combina un encoder testuale per la codifica vettoriale, un
vector store per la ricerca per similarità e un modello generativo; varianti successive
introducono re-ranking \citep{nogueira2020passage}, chunking avanzato e approcci
\emph{agentic RAG} in cui il retrieval diventa una decisione dinamica del modello
\citep{shi2023replug}.

La granularità dei chunk è cruciale: chunk piccoli aumentano la precisione semantica ma
perdono contesto, chunk grandi riducono il rumore ma penalizzano la similarity search. La
strategia \emph{Parent-Child} \citep{langchain2023parentdoc} risolve il compromesso con una
gerarchia a due livelli --- chunk figli piccoli ($\approx 400$ caratteri) per il retrieval,
chunk padre più ampi ($\approx 2000$ caratteri) restituiti al modello generativo. Il
\emph{contextual retrieval} di Anthropic \citep{anthropic2024contextual} estende l'idea
arricchendo ogni chunk con un'intestazione generata da un LLM che ne descrive il documento
di origine, migliorando la disambiguazione tra frammenti semanticamente simili.

LangGraph \citep{langgraph2024} estende LangChain con pipeline definite come grafi di stato
espliciti (StateGraph), superando i limiti delle LCEL chain lineari: ogni nodo opera su uno
stato condiviso (\code{TypedDict}) con transizioni deterministiche o content-based, offrendo
controllo fine-grained su ordine di esecuzione, gestione errori, persistenza della history
via checkpointer e iniezione di nodi di sicurezza senza toccare la logica principale.

La valutazione di sistemi RAG coinvolge sia la qualità del retrieval sia quella della
generazione. Ragas \citep{es2023ragas} propone \emph{Faithfulness} (affermazioni supportate
dal contesto), \emph{ResponseRelevancy} (pertinenza alla domanda) e \emph{ContextPrecision/
Recall} rispetto a una reference di ground-truth; un \emph{LLM-as-judge} \citep{zheng2023judging}
completa il quadro per criteri soggettivi come la coerenza.

I chatbot universitari hanno ricevuto crescente attenzione come strumenti di supporto agli
studenti: i sistemi rule-based o a intent classification \citep{smutny2020chatbots,
bansal2019systematic} si sono rivelati insufficienti a coprire la variabilità delle domande
studentesche, mentre gli LLM con RAG rendono praticabile un approccio ``open-domain''
ristretto al dominio. \sys{} si colloca in questo filone con la specificità di trattare
contenuti istituzionali italiani, multi-dominio e privi di metadati temporali strutturati.

%% ============================================================
\section{Architettura del Sistema}
\label{sec:architettura}

\subsection{Panoramica}

\sys{} è organizzato in tre sottosistemi principali: la \textbf{pipeline di ingestion}
(raccolta, estrazione, arricchimento e indicizzazione dei contenuti), il \textbf{core RAG
agentico} (grafo LangGraph per la gestione della conversazione e del retrieval) e la
\textbf{UI Gradio} per l'interazione con l'utente finale.
Dal punto di vista evolutivo, l'architettura è passata da una singola chain LCEL
concentrata in \code{src/brain.py} a un grafo LangGraph modulare, con stato, nodi e
orchestrazione separati nei file \code{src/agent/state.py}, \code{src/agent/nodes.py}
e \code{src/agent/brain.py}. Questa separazione ha reso espliciti i nodi di sicurezza
e ha ridotto l'accoppiamento tra guardrail, retrieval e generazione.

\begin{figure}[htbp]
  \centering
  \begin{tikzpicture}[
    scale=0.70, every node/.style={transform shape},
    node distance=0.8cm and 1.2cm,
    box/.style={draw,rounded corners=4pt,minimum width=2.6cm,minimum height=0.75cm,
    font=\small\sffamily,align=center,fill=blue!8,draw=blue!60!black},
  % Fixed cylinder proportions with aspect=0.25 to prevent it from bloating
    storage/.style={draw,cylinder,shape border rotate=90,aspect=0.25,minimum width=2.2cm,
    minimum height=1.5cm,font=\small\sffamily,align=center,
    fill=orange!10,draw=orange!70},
    arrow/.style={-{Stealth[length=5pt]},thick},
    groupbox/.style={draw=gray,dashed,rounded corners=5pt,inner sep=10pt}
  ]
    % --- Ingestion lane ---
    \node[box] (crawler) {Crawler\\multi-dominio};
    \node[box,right=of crawler] (parser) {Parser\\trafilatura};
    \node[box,right=of parser] (enricher) {Enrichment\\context header};
    \node[box,right=of enricher] (db) {Parent-Child\\Indexing};
    \node[storage,below=0.5cm of db] (chroma) {Chroma +\\LocalFileStore};

    \draw[arrow] (crawler) -- (parser);
    \draw[arrow] (parser) -- (enricher);
    \draw[arrow] (enricher) -- (db);
    \draw[arrow] (db) -- (chroma);

    \begin{scope}[on background layer]
      \node[groupbox,fit=(crawler) (parser) (enricher) (db) (chroma),
        label={[font=\footnotesize\bfseries]above:Pipeline di Ingestion}] {};
    \end{scope}

    % --- RAG lane ---
    % Increased vertical distance to 5cm to leave room for the Chroma database
    \node[box,below=5cm of crawler] (input_g) {input\_guard};
    \node[box,right=of input_g] (scope_g) {scope\_guard};
    \node[box,right=of scope_g] (reset) {reset\_state};
    \node[box,right=of reset] (agent) {agent\\(LLM)};
    \node[box,right=of agent] (tools) {tools\\rewrite/retrieve};
    \node[box,below=0.5cm of agent] (out_g) {output\_guard};
    \node[box,left=of out_g] (force) {force\_answer};

    \draw[arrow] (input_g) -- (scope_g);
    \draw[arrow] (scope_g) -- (reset);
    \draw[arrow] (reset) -- (agent);

    % Added yshift to prevent the text from overlapping the line
    \draw[arrow] (agent) -- node[above,font=\tiny,yshift=2pt]{tool calls} (tools);
    \draw[arrow] (tools) to[bend left=25] (agent);
    \draw[arrow] (agent) -- (out_g);
    \draw[arrow] (agent) -- (force);
    \draw[arrow] (force) -- (out_g);

    % Orthogonal routing from tools to chroma database
    \draw[arrow, rounded corners] (tools.north) |- (chroma.east);

    \begin{scope}[on background layer]
      \node[groupbox,fit=(input_g) (scope_g) (reset) (agent) (tools) (out_g) (force),
        label={[font=\footnotesize\bfseries]below:Core RAG Agentico (LangGraph)}] {};
    \end{scope}

    % --- User / UI ---
    \node[box,left=1.2cm of input_g,fill=green!8,draw=green!60!black] (ui) {UI Gradio\\(streaming)};

    \draw[arrow] (ui) to[bend left=10] (input_g);

    % Dropped the return line lower to bypass the LangGraph label text cleanly
    \draw[arrow, rounded corners] (out_g.south) -- ++(0,-1.2) -| (ui.south);

  \end{tikzpicture}
  \caption{Architettura generale di \sys{}. La pipeline di ingestion (alto) popola il vector store; il core RAG agentico (basso) gestisce le interazioni conversazionali.}
  \label{fig:arch}
\end{figure}

\subsection{Stack Tecnologico}

\sys{} è costruito su LangChain (loader, splitter, retriever, embedding), LangGraph per
l'orchestrazione del grafo agentico e Chroma come vector store.

\textbf{Embedding e Reranking.} Due modalità via \code{EMBEDDING\_PROVIDER}: locale
(\texttt{multilingual-e5-base}) o cloud (\texttt{qwen/qwen3-embedding-8b});
\texttt{baai/bge-m3} è stato incluso tra le alternative. Il reranking è locale
(\texttt{BAAI/bge-reranker-v2-m3}) o cloud (\texttt{Cohere Rerank~4 Pro}).

\textbf{Modelli LLM.} Due modalità via \code{LLM\_PROVIDER}: locale via Ollama
(\code{qwen2.5:7b} per agente e task lightweight) o cloud via OpenRouter
(\code{gpt-oss-120b} per l'agente, \code{llama-3.1-8b-instruct} per guardrail/rewrite),
quest'ultima la configurazione primaria. \code{gpt-oss-120b} è un modello open-weight a
117B parametri con architettura \emph{Mixture-of-Experts}: attiva solo 5,1B parametri
per forward pass, girando su una singola GPU H100 con quantizzazione MXFP4 nativa, e
supporta profondità di reasoning interno configurabile e tool use
nativo (function calling, structured output) \citep{openrouter2026gptoss120b} --- la combinazione di potenza di un
modello di larga scala e latenza da modello molto più piccolo ha motivato la scelta come
agente unico. La selezione del modello agente è stata comunque iterativa: da
\code{qwen2.5:7b} locale, a una fase intermedia \emph{planner-executor} (32B routing
$+$ 9B risposta), fino al singolo \code{gpt-oss-120b} finale in configurazione
\emph{single-agent} (il modello lightweight resta comunque separato per rewrite/guardrail,
mai per il routing dei tool), più stabile nel ciclo rewrite--retrieve--answer.

\textbf{UI.} Gradio~6 (\code{gr.ChatInterface}) con streaming token-per-token; l'indice
Chroma è caricato lazy all'avvio o ricostruito con il flag \code{--reindex}.

%% ============================================================
\section{Pipeline di Ingestion}
\label{sec:ingestion}

La pipeline di ingestion converte le pagine web del sito DIEM in chunk testuali
arricchiti, pronti per l'indicizzazione vettoriale. È composta da quattro fasi sequenziali:
crawling, parsing/filtraggio, arricchimento contestuale e indicizzazione. La traccia del
progetto definiva tre domini come perimetro della knowledge base:
\texttt{www.diem.unisa.it} (informazioni istituzionali, strutture, ricerca),
\texttt{docenti.unisa.it} (profili, pubblicazioni, ricevimento) e \texttt{corsi.unisa.it}
(programmi, materiale didattico). L'implementazione ha quindi dovuto gestire strutture di
navigazione e CMS eterogenei con strategie di crawling differenziate.

\subsection{Web Crawling Multi-Strategia}
\label{sec:crawling}

Il dominio principale \texttt{diem.unisa.it} viene analizzato tramite la sitemap HTML
(\texttt{/?sitemap}), che individua deterministicamente i seed principali, esplorati poi con
\code{RecursiveUrlLoader} (\code{max\_depth=3}); questo sostituisce un crawling ricorsivo
puro, non deterministico e incapace di raggiungere pagine con query string (es.\ dettagli
laboratori). I link verso \texttt{corsi.unisa.it} sono estratti dall'HTML grezzo
(\code{extract\_corsi\_urls}) e i PDF allegati scaricati con \code{load\_pdfs\_from\_links}.

Per i docenti, il sito DIEM non linka direttamente \texttt{docenti.unisa.it} ma rimanda a
schede \texttt{rubrica.unisa.it} con parametro matricola: la funzione
\code{extract\_diem\_faculty\_urls} risolve la matricola, verifica l'esistenza del link verso
\texttt{docenti.unisa.it} e scarta i profili privi di pagina dedicata, evitando di includere
personale di altri dipartimenti. La validazione è parallelizzata con un
\code{ThreadPoolExecutor} (10 worker) sfruttando il rate limiter globale; ogni docente
validato è poi crawlato separatamente (\code{max\_depth=2}) confinato al proprio profilo.
Analogamente, ogni corso è crawlato individualmente (\code{max\_depth=2}), limitando
\code{extract\_corsi\_urls} alla sezione \texttt{/offerta-formativa} per evitare corsi
obsoleti duplicati.

In tutte le fasi, \code{exclude\_dirs} esclude a monte risorse statiche e la funzione
\code{filter\_docs} rimuove URL con pattern indesiderati (\code{SKIP\_DOCUMENT\_SUBSTRINGS}),
applica deduplica per URL canonico e hash del contenuto ($\ge 350$ caratteri) e scarta
pagine in inglese/cinese. Il crawler usa una \code{RateLimitedSession} Token Bucket
(5 richieste/secondo) con retry automatico (\code{urllib3.util.retry}); il
\code{CrawlStateManager} mantiene un database SQLite di \texttt{ETag}/\texttt{Last-Modified}
per saltare, tramite HTTP~304, il ri-download di contenuti invariati nelle re-ingestion.

\subsection{Aggiornamento Incrementale Settimanale}
\label{sec:incremental-update}

Per evitare di ricostruire l'intero Chroma a ogni aggiornamento, il progetto prevede una
modalità incrementale orientata all'uso settimanale. Il principio è trattare ogni URL/source
come unità atomica di indicizzazione: in \code{db/crawl\_state.db} vengono salvati l'hash
del contenuto già indicizzato e i parent document ID associati. Durante un nuovo crawl, la
pipeline applica gli stessi passi della full ingestion (parsing, pulizia, deduplica e
filtri), raggruppa i documenti per \code{metadata["source"]} e ricalcola l'hash sul testo
post-filtro. Se l'hash coincide con quello persistito, la source viene saltata; se è nuova o
modificata, i vecchi child chunk e parent document di quella sola source vengono rimossi e
sostituiti tramite il \code{DocumentIndexer}.

La scelta è concettualmente corretta perché il confronto avviene sul contenuto che entrerebbe
realmente nel retriever, non sull'HTML grezzo. La cancellazione è localizzata e usa sia i
parent ID salvati nello stato sia i metadati presenti in Chroma, riducendo il rischio di
lasciare chunk orfani. Per prudenza, però, l'incrementale non elimina automaticamente tutte
le source che non ricompaiono nel crawl: una mancata visita può dipendere da timeout o pagine
temporaneamente non raggiungibili. Restano da monitorare variazioni non semantiche del
markup, modifiche ai filtri che cambiano gli hash, interruzioni tra delete e reindex e
metadati corrotti. Prima di automatizzarlo completamente sono quindi utili lock di
esecuzione, snapshot di Chroma/SQLite, dry-run con report delle source cambiate, verifica
post-run di coerenza child--parent--state e un full refresh periodico per rimuovere URL
realmente scomparsi.

È importante distinguere aggiornamento e configurazione sperimentale: il Chroma usato nella
configurazione finale più robusta non coincide necessariamente con l'ultimo indice prodotto
dall'update incrementale. Anche a parità di parametri, context header generati da modello,
ordine del crawling e piccole variazioni di parsing possono produrre store leggermente
diversi; l'ultimo indice è quindi il più aggiornato, mentre quello selezionato per C5 è il
più robusto rispetto alle valutazioni effettuate.

La pipeline integra fonti supplementari: il modulo \code{easycourse.py} interroga le API
di EasyCourse per calendario esami e orario lezioni dei cinque corsi DIEM, producendo un
documento per coppia (CDL, insegnamento). Le risposte JSON vengono normalizzate e
raggruppate per insegnamento, così un singolo retrieve restituisce appelli o lezioni
rilevanti senza frammentazione. Vengono inoltre generati gli URL dei bandi DIEM (8 categorie) e scaricate le sottopagine
\texttt{/focus?id=\ldots} non raggiungibili con \code{max\_depth=1}; una sorgente statica
reinietta dati amministrativi (P.IVA, sede) rimossi dal parser come boilerplate ma
richiesti da domande atomiche.

\subsection{Parsing e Filtraggio Temporale}
\label{sec:parser}

Il testo è estratto con \textbf{trafilatura} \citep{barbaresi2021trafilatura} (filtro
\code{favor\_precision=True}). I PDF usano
\texttt{pdfplumber}, con le pagine dello stesso documento unite prima dello split per
abilitare chunk cross-page su regolamenti multi-pagina.

Trafilatura classifica erroneamente come ``non-content'' le pagine con struttura
accordion/tabellare (pannelli corti), producendo output nulli. Per un sottoinsieme di URL
critici (pagine personale, organi collegiali, commissioni, contatti corso, didattica
docente) la funzione \code{html\_extractor\_for\_source} attiva un'estrazione BeautifulSoup
mirata via whitelist (\code{\_is\_structured\_source}); per tutti gli altri URL resta
trafilatura primario con fallback BeautifulSoup generico. Un parser BeautifulSoup globale è
stato scartato perché la funzione di rimozione boilerplate troncava le biografie dei
docenti: la whitelist mirata contiene questo rischio.

Poiché i CMS non espongono sempre metadati temporali strutturati, il filtro
\code{filter\_recent\_documents} applica un cutoff minimo al 2020: se un documento contiene
un anno esplicito in URL, metadati o testo analizzato, e l'anno più recente trovato è
$<2020$, il documento viene scartato; in assenza di segnali temporali espliciti viene
mantenuto per evitare falsi negativi su pagine istituzionali stabili. Il crawling introduce
inoltre vincoli mirati a monte: gli URL dei bandi DIEM sono generati solo per anno corrente
e anno precedente (nella run 2026: 2026 e 2025), mentre le pubblicazioni dei docenti sono
crawlate esplicitamente per tutti gli anni dal 2020 all'anno corrente. I documenti scartati
sono salvati in \texttt{dropped\_temporal\_filter.json} per audit.

\subsection{Arricchimento Contestuale: Context Header}
\label{sec:enrichment}

La pipeline di arricchimento contestuale risolve un problema strutturale del retrieval su
documenti universitari: frammenti come ``modalità d'esame'', ``ricevimento'' o ``obiettivi
formativi'' sono semanticamente identici tra decine di pagine diverse e, senza contesto
negli embedding, generano falsi positivi sistematici. La qualità dell'header è garantita
da un routing ibrido \emph{per-URL} in \code{enrichment.py}: per i domini con metadati
titolo affidabili (\texttt{docenti.unisa.it}, \texttt{corsi.unisa.it}, EasyCourse) si usa
l'euristico deterministico di \code{classify\_context\_header}; per i PDF di
\texttt{corsi.unisa.it} e per le pagine di \texttt{diem.unisa.it} si invoca l'LLM
(\code{mistralai/mistral-nemo} via OpenRouter, con fallback Ollama ed euristico), reso
configurabile tramite \code{OPENROUTER\_CONTEXT\_HEADER\_MODEL} per usare
\code{llama-3.1-8b-instruct} nelle run finali senza modifiche al codice. La funzione
\code{repair\_context\_header\_semantics} corregge post-hoc classificazioni errate su
URL-pattern noti; l'header è limitato a 18 parole, include un year tag
\texttt{[AY YYYY/YYYY]} quando rilevabile e viene cachato su disco.

L'header non viene inserito nel corpo del documento ma salvato nei metadati come
\code{context\_header}: in \code{database.py}, dopo la generazione dei chunk padre, lo
splitter \code{ContextHeaderTextSplitter} lo prepende a ogni chunk figlio prima
dell'indicizzazione, mentre il parent resta pulito nel \code{LocalFileStore}. Questo
corregge il limite della prima versione, dove l'header nel testo intero pre-split veniva
ereditato solo dal primo chunk figlio; la validazione su seed stratificati
(\code{simulate\_enrichment.py}) ha confermato un miglioramento della qualità degli header
del 93--95\%. Per ridurre costo e incoerenza, un grouping per sorgente genera un solo
header per URL con meno di dieci parent chunk (assegnato a tutti i chunk della pagina),
mantenendo l'header per-chunk sui documenti più grandi dove sezioni diverse trattano
argomenti distinti.

\begin{tcolorbox}[boxstyle,title={Esempio di context header applicato a un chunk figlio}]
\textbf{Context header:} \textit{Corso di Ingegneria del Software -- prof. Mario Rossi -- orari esame a.a. 2025/2026}\\[4pt]
\textbf{Testo del chunk:} \textit{Modalità d'esame: prova scritta (60\%) + orale (40\%). La prova scritta si svolge in laboratorio...}
\end{tcolorbox}

\subsection{Strategia di Chunking Parent-Child}
\label{sec:chunking}

La struttura Parent-Child è implementata con i seguenti parametri (Tabella~\ref{tab:chunking}),
scelti empiricamente per bilanciare precisione semantica dei chunk figli e ricchezza
contestuale dei chunk padre. La taratura non è stata monotona: sono state provate sia
configurazioni più piccole, per ridurre il rumore, sia configurazioni più grandi, per
contenere schede corso e formule di regolamento in un singolo parent. La configurazione
finale C5 torna a 2000/200 e 400/50 perché l'\emph{adjacent retrieval} recupera il chunk
successivo quando serve continuità, evitando di aumentare sempre la dimensione del
contesto.

\begin{table}[htbp]
\centering
\caption{Principali configurazioni di chunking valutate. Le dimensioni sono in caratteri.}
\label{tab:chunking}
\small
\begin{tabularx}{\textwidth}{lrrX}
\toprule
\textbf{Fase} & \textbf{Parent} & \textbf{Child} & \textbf{Motivazione/esito} \\
\midrule
Baseline & 2000/200 & 400/50 & Buon equilibrio iniziale, ma alcune formule PDF e schede corso restavano frammentate. \\
v2.1 & 1500/180 & 350/70 & Chunk più densi e meno rumorosi; miglioramento su pagine list-heavy, ma perdita di contesto su regolamenti lunghi. \\
v3 / C4 & 3750/375 & 600/80 & Copre schede corso e formule in parent più autocontenuti; aumenta però il rumore e la lunghezza del contesto. \\
C5 finale & 2000/200 & 400/50 & Ritorno a chunk più compatti, compensato da adjacent retrieval post-rerank. \\
\bottomrule
\end{tabularx}
\end{table}

Il \code{ParentDocumentRetriever} di LangChain viene configurato per usare il vector store
Chroma (chunk figli) per la ricerca per similarità e il \code{LocalFileStore} su disco
per il recupero dei chunk padre corrispondenti. Il retrieval restituisce quindi il chunk
padre, più ricco di contesto, come input al modello generativo. L'indicizzazione avviene
in batch di massimo 100 chunk figli per volta per gestire la memoria durante l'inserimento
in Chroma.

\subsection{Embedding e Indicizzazione}

Il sistema integra la classe \code{OpenRouterEmbeddings} per gli embedding cloud. La
versione finale usa \texttt{qwen/qwen3-embedding-8b}: il wrapper applica alle query il
prefisso \texttt{Instruct: Retrieve relevant passages that answer the user's query} seguito
da \texttt{Query: \ldots}, mentre i documenti sono indicizzati senza prefisso, seguendo il
comportamento instruction-tuned del modello. In modalità locale, \code{E5HuggingFaceEmbeddings}
mantiene la codifica simmetrica di E5 (\texttt{query:}/\texttt{passage:}); \texttt{e5-base}
è stato usato come baseline locale perché multilingue, leggero e riproducibile senza API
esterne, con embedding a 768 dimensioni gestibili in Chroma. L'evoluzione
dell'embedding ha attraversato \texttt{multilingual-e5-small/base} (locale) e
\texttt{baai/bge-m3} (cloud) prima di convergere su Qwen3. BGE-M3 è stato scartato non per
scarsa qualità ma perché il progetto usa Chroma dense-only (senza BM25/ColBERT che ne
sfruttino la natura ibrida); Qwen3, pur a 4096 dimensioni, ha dato il miglior recall nelle
configurazioni C3--C5.

La \code{similarity\_score\_threshold} è impostata a 0,50 con $k=15$ candidati iniziali,
poi re-scored dal reranker: rispetto a un setup locale iniziale (\texttt{e5-base}, soglia
0,70, $k=20$), la soglia è stata abbassata perché 0,70 escludeva documenti rilevanti su
query brevi, mentre $k$ è stato ridotto a 15 dopo aver verificato che i candidati oltre
quella posizione venivano quasi sempre scartati dal reranker.

%% ============================================================
\section{Core RAG Agentico}
\label{sec:core-rag}

\subsection{Evoluzione dell'Architettura: da Chain a Grafo}

L'architettura del core RAG ha attraversato quattro fasi evolutive: una \emph{chain LCEL}
a flusso fisso (rewrite $\to$ retrieve $\to$ rerank $\to$ format $\to$ LLM, nessun tool
calling); un \emph{agent + MemorySaver} con tool \code{retrieve}/\code{answer}, dove
l'agente non chiamava sempre retrieve (rischio allucinazione); uno \emph{StateGraph
esplicito} con nodi separati e un \code{retrieve\_node} obbligatorio su due modelli
(32B+9B); infine l'architettura \emph{fully agentic} finale, dove il retrieve\_node fisso
è rimosso e l'agente decide autonomamente quando recuperare, con un solo modello.

La transizione chiave è l'ultima: rimuovere il nodo \code{retrieve\_node} obbligatorio
aumenta la flessibilità al costo di una LLM call aggiuntiva per turno, compensata da un
hint dinamico iniettato nel system prompt quando il contatore di retrieve è zero, che
forza il modello a chiamare \code{retrieve} come primo tool del turno senza fallback in
codice.

\subsection{LangGraph StateGraph: Struttura del Grafo}
\label{sec:graph}

Il grafo è definito come un \code{StateGraph} LangGraph con stato condiviso di tipo
\code{DiemState}:

\begin{lstlisting}[language=Python, caption={Definizione di DiemState (src/agent/state.py)}]
class DiemState(TypedDict):
    messages: Annotated[List[BaseMessage], add_messages]
    tool_call_count: int    # retrieve calls only
    retrieved_context: str  # last retrieved text
    last_docs: List[Document]  # docs from last retrieve
\end{lstlisting}

I campi \code{tool\_call\_count}, \code{retrieved\_context} e \code{last\_docs} vengono azzerati
all'inizio di ogni turno dal nodo \code{reset\_state}, garantendo che il contatore e il
contesto non si propaghino tra domande diverse.

Il grafo comprende otto nodi con le funzioni descritte nella Tabella~\ref{tab:nodes}.

\begin{table}[htbp]
\centering
\caption{Nodi del LangGraph StateGraph di \sys{}.}
\label{tab:nodes}
\small
\begin{tabularx}{\textwidth}{lX}
\toprule
\textbf{Nodo} & \textbf{Funzione} \\
\midrule
\code{input\_guard}    & Controllo contenuto offensivo sull'input utente (regex + badwords). \\
\code{scope\_guard}    & Classificatore di ambito: DIEM/UNISA vs.\ fuori dominio. \\
\code{reset\_state}    & Azzera contatori e contesto per il nuovo turno. \\
\code{agent}           & Loop principale: invoca l'LLM con binding dei tool. \\
\code{tools}           & Esegue i tool richiesti dall'agente; incrementa retrieve counter. \\
\code{forced\_retrieve}& Safety net: inietta retrieve sintetico se l'agente ha generato output vuoto senza retrieve. \\
\code{force\_answer}   & Stub i tool calls pendenti se raggiunto MAX\_RETRIEVE\_CALLS; forza risposta finale. \\
\code{output\_guard}   & Controllo offensivo + redazione PII sull'output del sistema. \\
\bottomrule
\end{tabularx}
\end{table}

Il flusso principale è:
\begin{multline*}
\text{START} \to \text{input\_guard} \to \text{scope\_guard} \to \text{reset\_state} \\
\to \text{agent} \leftrightarrow \text{tools} \to \text{output\_guard} \to \text{END}
\end{multline*}

\noindent Le deviazioni rispetto al flusso principale gestiscono i casi di rifiuto
(\code{input\_guard} o \code{scope\_guard} iniettano un \code{AIMessage} e terminano),
di loop forzato (\code{forced\_retrieve} e \code{force\_answer}) e di esaurimento
del budget di retrieve (\code{agent} $\to$ \code{force\_answer} $\to$ \code{output\_guard}).

\subsection{Guardrail a Cascata}
\label{sec:guardrail}

Il sistema implementa tre livelli di guardrail indipendenti.

\textbf{Controllo contenuto offensivo.} La classe \code{OffensiveContentGuardrail} mantiene
una lista di parole inappropriate (250+ termini, codificata in base64) e applica pattern
matching regex sull'input (\code{input\_guard}) e sull'output (\code{output\_guard}). Il
controllo è puramente regex: su un dominio testuale ristretto come quello universitario la
lista è stabile e i falsi negativi trascurabili, rendendo superfluo un componente LLM
dedicato.

\textbf{Scope guardrail.} La classe \code{ScopeGuardrail} implementa un classificatore
gerarchico: un fast-path su circa 50 keyword del dominio DIEM/UNISA considera in-scope la
domanda senza invocare l'LLM; in assenza di keyword, un modello lightweight applica un
prompt binario che descrive esplicitamente le categorie da bloccare (sport, cucina,
personaggi storici non accademici) e quelle da lasciar passare (persone in contesto
accademico, topic ambigui). Questo disegno riduce i falsi negativi osservati su domande
tipo ``Chi è [nome professore]?'', inizialmente bloccate per assenza di keyword esplicite.
Per i follow-up multi-turno, se esiste history precedente il nodo riscrive prima la
domanda con il modello lightweight e valuta lo scope sulla query risolta, evitando che
messaggi brevi come ``Davvero?'' vengano bloccati solo perché privi di keyword DIEM se
isolati dal contesto.

\textbf{Redazione PII.} La funzione \code{redact\_pii} cerca nell'output pattern di
indirizzi e-mail e numeri di carta di credito (quattro blocchi separati da spazi/trattini,
per evitare falsi positivi su codici identificativi generici), sostituendo il match con un
messaggio standard di privacy.

\subsection{Gestione della Storia Conversazionale}

La storia conversazionale è gestita da \code{MemorySaver} (LangGraph) con chiave
\code{session\_id}. Prima di passare i messaggi all'agente, il nodo \code{agent} applica
una fase di pulizia che rimuove i \code{ToolMessage} dei turni precedenti, rimuove gli
\code{AIMessage} con tool calls pendenti già risolti, e sostituisce gli \code{AIMessage}
iniettati dai guardrail con un placeholder neutro, così che il modello veda coppie
Human/AI semanticamente pulite. Questo riduce il token bloat rispetto alle versioni
iniziali, dove l'accumulo di ToolMessage con i documenti recuperati saturava rapidamente
il context window.

\textbf{Gestione delle sessioni lunghe: da \texttt{summarize\_history} a sliding window.}
Per le sessioni oltre i dieci turni, un primo approccio ha introdotto un tool
\code{summarize\_history} per comprimere la history in bullet list al decimo turno. I test
hanno rivelato un problema di compliance: dopo la chiamata il modello generava una
risposta generica senza procedere con \code{retrieve()}, interpretando il summary come
contesto sufficiente --- la sequenza multi-tool nello stesso turno risultava inaffidabile
per questo modello, e il tool è stato rimosso.

La soluzione adottata è una \textbf{sliding window deterministica}: il nodo \code{agent}
mantiene intatti gli ultimi 5 turni e rimuove \code{ToolMessage} e \code{AIMessage} con
tool calls pendenti dai turni più vecchi, prima di costruire la sequenza passata al
modello. Il \code{MemorySaver} conserva lo stato completo del grafo; la window agisce solo
sulla vista ricevuta dal modello, senza round trip LLM aggiuntivo. Il sistema registra in
log il numero di messaggi rimossi
(\texttt{agent | sliding\_window | kept\_turns=5 | dropped\_tool\_msgs=N}).

\begin{lstlisting}[language=Python, caption={Sliding window in \texttt{\_node\_agent} (src/agent/nodes.py): strip dei ToolMessage dai turni fuori finestra.}]
human_indices = [i for i, m in enumerate(state["messages"])
                 if isinstance(m, HumanMessage)]
cutoff_idx = human_indices[-5] if len(human_indices) > 5 else 0

clean_messages = []
for i, m in enumerate(state["messages"]):
    # strip tool traffic from turns outside the 5-turn window
    if i < cutoff_idx and (
        isinstance(m, ToolMessage)
        or (isinstance(m, AIMessage) and getattr(m, "tool_calls", None))
    ):
        continue
    clean_messages.append(m)
\end{lstlisting}

\subsection{Pipeline di Retrieval}
\label{sec:retrieval}

La pipeline di retrieval segue uno schema a due stadi comunemente chiamato
\emph{retrieve and rerank}.

\textbf{Primo stadio --- Bi-encoder.} Il \code{ParentDocumentRetriever} LangChain cerca per
similarità coseno nel vector store Chroma sui chunk figli, con $k=15$ candidati e soglia
0,50 (\code{RETRIEVER\_SCORE\_THRESHOLD}). Il
bi-encoder finale, come anticipato nella Sezione~\ref{sec:ingestion}, è
\texttt{qwen/qwen3-embedding-8b} via OpenRouter.

\textbf{Secondo stadio --- Reranker.} I 15 chunk figli vengono re-scored da un secondo
modello sulla coppia (query, chunk). Il reranker è dual-provider: localmente il
cross-encoder \code{BAAI/bge-reranker-v2-m3} (singleton a livello di modulo, per eliminare
un bottleneck di latenza $\approx 20$s osservato inizialmente); in cloud
\code{Cohere Rerank~4 Pro} (32K contesto) via OpenRouter, adottato per maggiore finestra
di contesto e robustezza multilingue. I cinque chunk con score più alto vengono
selezionati e per ciascuno viene recuperato il chunk padre corrispondente dal
\code{LocalFileStore}.

\textbf{Adjacent retrieval post-rerank.} Per mitigare la frammentazione residua dei parent,
ogni chunk padre riceve tre metadati (\code{chunk\_index}, \code{chunk\_total},
\code{chunk\_id} hash di \texttt{source:chunk\_index}); dopo il reranking, il tool
\code{retrieve} recupera anche il chunk immediatamente successivo per ciascun parent
selezionato tramite \code{mget}. Lo spostamento \emph{dopo} il reranking è deliberato: il
chunk adiacente è una continuazione di un chunk già giudicato rilevante e non deve
competere isolatamente nel reranker, che altrimenti potrebbe scartarlo perché privo della
frase iniziale.

\textbf{Recency boost.} Dopo il reranking viene applicato un piccolo bonus di rilevanza
basato sull'anno accademico estratto dal context header, esclusivamente sui PDF, con
entità trascurabile: agisce da tiebreaker tra documenti con score quasi identico,
privilegiando la versione più recente senza alterare l'ordinamento generale.

\begin{lstlisting}[language=Python, caption={Recency boost applicato dopo il reranking (\texttt{reranker.py}).}]
_YEAR_RE = re.compile(r'\[AY (\d{4})|\[(\d{4})\]')

def _apply_recency_boost(docs):
    for doc in docs:
        year = _extract_year(doc.metadata.get("context_header", ""))
        is_pdf = doc.metadata.get("source", "").lower().endswith(".pdf")
        bonus = max(0.0, (year - 2020) * 0.001) if (year and is_pdf) else 0.0
        base = doc.metadata.get("relevance_score", 0.0)
        doc.metadata["relevance_score_boosted"] = round(base + bonus, 6)
    return sorted(docs, key=lambda d: d.metadata["relevance_score_boosted"],
                  reverse=True)
\end{lstlisting}

L'anno viene estratto tramite espressione regolare dai tag \texttt{[AY 2025/2026]}
o \texttt{[2021]} già presenti nel context header: nessuna chiamata aggiuntiva al
modello è necessaria. Il bonus massimo è $0{,}005$ per anno 2025, trascurabile
rispetto agli score del reranker ($\in [0,1]$).

\subsection{Tool dell'Agente}
\label{sec:tools}

L'agente dispone di due tool LangGraph, descritti nella Tabella~\ref{tab:tools}. I tool
\code{summarize} e \code{calculate}, presenti nelle versioni precedenti, sono stati
rimossi nel corso dello sviluppo per ragioni di affidabilità e qualità descritte di
seguito.

\begin{table}[htbp]
\centering
\caption{Tool disponibili all'agente di \sys{}.}
\label{tab:tools}
\small
\begin{tabularx}{\textwidth}{lX}
\toprule
\textbf{Tool} & \textbf{Descrizione e vincoli operativi} \\
\midrule
\code{rewrite(query)} & Riscrive la query in una domanda standalone risolvendo
                        pronomi, riferimenti impliciti e query incomplete. Per contenuto
                        accademico senza anno appende automaticamente
                        ``a.a.\ 2025/2026''. Chiamato obbligatoriamente prima di
                        \emph{ogni} \code{retrieve}, inclusi i retry: al secondo
                        tentativo produce automaticamente una query diversificata
                        rispetto al precedente output. Non soggetto a cap. \\
\addlinespace
\code{retrieve(query)} & Esegue il retrieval bi-encoder + reranker sulla knowledge base
                         DIEM. Aggiorna \code{retrieved\_context} e \code{last\_docs} nello stato.
                         Obbligatorio almeno una volta per turno; soggetto a cap
                         \code{MAX\_RETRIEVE\_CALLS=3} (solo questo tool conta verso il limite). \\
\bottomrule
\end{tabularx}
\end{table}

Solo \code{retrieve} incrementa \code{tool\_call\_count}. Questo design deliberato
permette all'agente di chiamare \code{rewrite} liberamente senza avvicinarsi al
limite di retrieve, che è il collo di bottiglia computazionale.

\textbf{Rimozione di \texttt{summarize} e \texttt{calculate}.} Il tool
\code{summarize(text, query)}, pensato per condensare contesti retrieved lunghi, non
veniva quasi mai chiamato dal 120b, che gestiva direttamente contesti estesi. La variante
\code{summarize\_history} è stata rimossa per le ragioni discusse nella gestione della
storia conversazionale; entrambe le forme di sintesi sono state sostituite da una gestione
deterministica della finestra di contesto.

\newpage
Il tool \code{calculate(ctx, op, vals)} eseguiva calcoli accademici (medie ponderate, voto
di laurea) delegando al modello lightweight tramite \texttt{simpleeval}. Un benchmark su
tre scenari reali (Tabella~\ref{tab:calc_bench}) ha evidenziato due problemi: sulla
correttezza, \texttt{simpleeval} crashava su sintassi tuple generate dall'8b (es.\
\texttt{min(((PT+PC),5),5)} $\to$ \textit{``Tuple is not available in this evaluator''}),
con un secondo tentativo che produceva spesso risultati errati; sulla latenza, l'overhead
del tool portava il TTFT da 15--33\,s contro i 7--10\,s senza tool.

\begin{lstlisting}[language=Python, caption={Schema della prima implementazione del tool \texttt{calculate}: valutazione controllata di espressioni aritmetiche generate dal modello lightweight.}]
@tool
def calculate(ctx: str, op: str, vals: dict) -> str:
    """Compute academic formulas such as weighted average or graduation grade."""
    expr = build_expression_from_context(ctx, op, vals)
    result = simple_eval(expr, names=vals)
    return json.dumps({"expression": expr, "result": result})
\end{lstlisting}

Lo snippet sintetizza la versione rimossa: il modello produceva un'espressione simbolica
valutata in modo controllato dal tool, isolando l'aritmetica dal modello generativo. Nella
pratica, la qualità dell'espressione dipendeva ancora dal modello lightweight; quando
malformata, il tool diventava un punto di fallimento invece che una garanzia.

\begin{table}[htbp]
\centering
\caption{Benchmark \texttt{calculate} vs.\ inline per tre scenari di voto di laurea
         (media 27, TTFT in secondi).}
\label{tab:calc_bench}
\small
\begin{tabular}{lrrl}
\toprule
\textbf{Scenario} & \textbf{TTFT con tool} & \textbf{TTFT inline} & \textbf{Qualità} \\
\midrule
Magistrale IngInf (IE227) & 15,2\,s & 10,3\,s ($-32\%$) & equivalente \\
Triennale IngInf (IE127)  & 33,0\,s & 10,2\,s ($-69\%$) & \textbf{inline superiore} \\
Magistrale IEDM (IE232)   & 18,5\,s & \phantom{0}7,4\,s ($-60\%$) & equivalente \\
\bottomrule
\end{tabular}
\end{table}

Il caso triennale è stato decisivo: il tool crashava sulla sintassi tuple, il secondo
tentativo produceva un valore errato con variabili nulle, mentre \code{gpt-oss-120b}
inline restituiva la formula completa con tutti i parametri (cap, range PT, costante FC,
divisore) e una tabella di scenari al variare di PT. La rimozione del tool è quindi
motivata non da ragioni di semplicità, ma da affidabilità superiore e latenza ridotta.

\textbf{Rewrite retry-aware.} L'architettura agentica permette più retrieve nello stesso
turno finché il contesto non è sufficiente, ma questa capacità risultava inutile in
pratica: al retry, \code{rewrite} produceva una query quasi identica alla precedente,
recuperando gli stessi documenti. Il tool è stato quindi esteso per rilevare il retry
(controllando se esiste già un \code{ToolMessage} di \code{rewrite} nel turno corrente) e,
in caso affermativo, generare una query semanticamente diversa: se la query precedente
aveva aggiunto termini non richiesti (es.\ ``DIEM''), il retry si avvicina alla
formulazione originale dell'utente; altrimenti cambia livello di specificità, sinonimi o
forma interrogativa. Un vincolo esplicito nel prompt (\texttt{CRITICAL: your output MUST
differ from <previous\_query>}) forza la diversificazione, e \code{rewrite()} non è mai
considerato azione terminale: ogni chiamata deve essere seguita da una \code{retrieve},
correzione introdotta dopo aver osservato casi in cui l'agente riscriveva la domanda e poi
rispondeva senza recuperare contesto.

\subsection{Prompt Engineering}
\label{sec:prompts}

La progettazione dei prompt ha richiesto iterazione significativa, poiché piccole
variazioni producono comportamenti qualitativamente diversi.
Il sistema combina quindi più tecniche di prompting (role prompting, structured prompting
via XML, few-shot e meta-prompting), ciascuna introdotta per mitigare un problema
osservato empiricamente.

\textbf{AGENT\_SYSTEM\_PROMPT.} Il prompt è prevalentemente \emph{zero-shot}: usa
\emph{role prompting}, step procedurali (\emph{tool use}, \emph{context validation},
\emph{generate answer}) e tag strutturati \texttt{[KNOWLEDGE\_GAP]}/\texttt{[FUORI\_SCOPE]},
con esempi limitati a casi di ambiguità o date. Le regole operative, incluse quelle su
date, ruoli accademici e calcoli multi-step, restano istruzioni nude: per il comportamento
agentico procedurale è sufficiente il meta-prompting, senza esempi numerici specifici.
Alcune regole derivano da errori osservati: ``NEVER EMPTY'' per evitare soli tag senza
spiegazione, soglia knowledge-gap più morbida su contesto parziale e regola sui ruoli
accademici dopo downgrade errati di docenti PO/PA a ricercatori.

\textbf{REWRITE\_PROMPT.} Il rewrite è invece \emph{few-shot}: regole sui casi edge
(pronomi, entità da preservare, anno accademico, challenge di conferma) sono seguite da
esempi input/output annotati, perché la riscrittura è un task più ambiguo e pattern-based.
I casi su voto di laurea, aree di ricerca, laboratori e piano di studi sono stati aggiunti
empiricamente. Tag XML-like in \code{AGENT\_SYSTEM\_PROMPT}/\code{REWRITE\_PROMPT}
strutturano i few-shot examples e separano dati da istruzioni.
La regola 7 (voto di laurea) applica \emph{step-back prompting} \citep{zheng2023stepback}:
la query specifica (``con media 27...'') è riscritta in una domanda astratta (``qual è la
formula...''), il retrieve recupera il regolamento generale, e l'agente ricompone il
risultato specifico nello Step~3.

\textbf{Temperature.} Il modello agente usa temperatura 0,1 per mantenere una minima
flessibilità nel reasoning su task fattuali; il modello lightweight usa temperatura 0,
rendendo deterministici guardrail, classificazione e riscrittura.

\textbf{Gestione del voto di laurea.} Il calcolo richiede parametri specifici per corso
(formula FC, costanti, bonus, differenti tra i cinque corsi DIEM), gestiti con tre
meccanismi: un \emph{cancello di chiarimento} allo Step~1 (se il corso non è specificato,
l'agente chiede chiarimento prima di ogni tool, per evitare un calcolo su regolamento
sbagliato); un \emph{controllo del regolamento} allo Step~2 (verifica che l'URL del
documento recuperato contenga il codice del corso richiesto, altrimenti nuovo retrieve
mirato). La rimozione del tool \code{calculate} non elimina il supporto ai calcoli
accademici: il calcolo passa all'agente principale, guidato da \emph{meta-prompting}
\citep{suzgun2024metaprompting} a passaggi espliciti. Per il voto di laurea, il prompt
impone di identificare prima media, FC, bonus e vincoli del regolamento, e solo dopo
produrre il risultato finale.

\textbf{Tag di rifiuto.} I tag \texttt{[FUORI\_SCOPE]} e \texttt{[KNOWLEDGE\_GAP]} sono
marcatori strutturati nell'output; un dizionario \code{\_TAG\_FALLBACKS} in
\code{brain.py} mappa ciascun tag a un messaggio in italiano quando è l'unico contenuto
dell'output.

%% ============================================================
\section{Framework di Valutazione}
\label{sec:evaluation}

\subsection{Golden Set}

Il dataset di valutazione (\emph{golden set}) è composto da domande suddivise in quattro
categorie progettate per coprire le modalità d'uso principali e i casi critici del sistema:

\begin{itemize}
  \item \textbf{in\_scope}: domande fattuali sul DIEM (corsi, docenti, strutture, procedure);
  \item \textbf{multi\_turn}: sequenze di domande in cui i turni successivi contengono
        pronomi o riferimenti impliciti al turno precedente;
  \item \textbf{out\_of\_scope}: domande deliberatamente fuori ambito (sport, cucina,
        personaggi storici) per misurare la capacità di rifiuto;
  \item \textbf{robustness}: scenari dedicati al criterio custom descritto nella
        sottosezione seguente.
\end{itemize}

La valutazione finale usa il golden set italiano \code{golden\_set\_it.json}: le domande
con \code{reference} esplicita alimentano le metriche Ragas ground-truth, le altre solo le
metriche no-reference. Il dataset è stato costruito manualmente con domande realistiche
sui contenuti DIEM; le categorie dedicate ai casi fuori ambito e alle challenge
conversazionali sono state aggiunte in un secondo momento, dopo i primi test manuali.

\subsection{Metriche di Valutazione}

Il framework utilizza Ragas~0.4.3 \citep{es2023ragas} con sei metriche (Tabella~\ref{tab:metrics}).

\begin{table}[htbp]
\centering
\caption{Metriche di valutazione del framework Ragas.}
\label{tab:metrics}
\small
\begin{tabularx}{\textwidth}{lXl}
\toprule
\textbf{Metrica} & \textbf{Descrizione} & \textbf{Reference?} \\
\midrule
ResponseRelevancy      & Misura la pertinenza della risposta alla domanda.          & No \\
Faithfulness           & Proporzione di affermazioni supportate dal contesto.       & No \\
AnswerCorrectness      & Correttezza fattuale rispetto alla reference.              & Sì \\
LLMContextPrecision    & Precisione del contesto recuperato rispetto alla reference.& Sì \\
LLMContextRecall       & Recall del contesto recuperato rispetto alla reference.    & Sì \\
Coherence (AspectCritic) & Coerenza stilistica e logica della risposta (LLM-judge).& No \\
\bottomrule
\end{tabularx}
\end{table}

\noindent Per evitare valori NaN artificiali su righe senza reference (che Ragas tratta come
errori di valutazione), il runner implementa una valutazione a due fasi: le metriche
no-reference vengono calcolate sull'intera batteria di domande, mentre le metriche
reference-requiring vengono calcolate esclusivamente sulle righe con campo \code{reference}
non vuoto.

\subsection{Criteri Custom}

I due criteri custom sono calcolati fuori da Ragas come pass-rate sulle categorie dedicate
del golden set, perché misurano comportamenti conversazionali e di sicurezza non coperti
dalle metriche standard.

\textbf{Scope Awareness.} Per ogni domanda out-of-scope, il runner verifica se la
risposta contiene un rifiuto esplicito (\emph{strict pass}: marker \texttt{[FUORI\_SCOPE]}
o frasi standard) o un rifiuto implicito per knowledge gap (\emph{soft pass}: il sistema
dichiara di non avere informazioni senza inventare contenuto), con un LLM-judge
lightweight sui casi incerti. La distinzione è metodologicamente importante: senza di
essa il tasso di scope awareness risulterebbe artificialmente vicino al 100\%,
mascherando i casi in cui il sistema non identifica attivamente la domanda come fuori
ambito.

\textbf{Robustness.} Per ogni scenario avversariale, il runner verifica che il sistema
non cambi risposta senza motivo dopo una challenge (``Sei sicuro?''), non accetti una
false premise e non esegua istruzioni di jailbreak, con override su marker strutturati
quando il comportamento atteso è il rifiuto esplicito.

\subsection{Output del Runner}

Per ogni run, il runner produce file strutturati (\code{run.log},
\code{per\_question.json}, metriche Ragas/custom e \code{summary.md}) e una cache
\code{off}/\code{use}/\code{refresh} indicizzata su schema, modello, temperatura,
sessione, history e domanda, evitando condivisioni tra configurazioni diverse.

%% ============================================================
\section{Validazione Intermedia e Decisioni di Stabilizzazione}
\label{sec:stabilization}

Le prime valutazioni hanno avuto soprattutto una funzione diagnostica: non servivano a
scegliere la configurazione finale, ma a individuare i punti in cui il comportamento del
sistema non era ancora stabile. Una smoke run iniziale con modello locale mostrava valori
di Faithfulness nell'intervallo 0,490--0,504, incompatibili con l'obiettivo di un
assistente universitario affidabile; il divario Scope Awareness strict/soft (33,3\% vs
100\%) indicava knowledge gap riconosciuti ma fuori ambito non marcati esplicitamente.
L'analisi manuale delle
risposte e dei log ha ricondotto il problema a quattro criticità principali.

\textbf{Retrieve non garantito e knowledge gap.} Nelle prime architetture agentiche il
modello poteva rispondere senza richiamare il retrieval, oppure usare un contesto parziale
come se fosse sufficiente. Per questo il grafo finale introduce un hint dinamico quando il
contatore di retrieve è zero, il safety net \code{forced\_retrieve}, tag strutturati
\code{[KNOWLEDGE\_GAP]}/\code{[FUORI\_SCOPE]} e regole di citazione più esplicite nel
prompt.

\textbf{Drift nei follow-up e rewrite generico.} Nei turni multi-turno, messaggi brevi
come ``Sei sicuro?'' o domande con pronomi potevano generare un nuovo retrieve su query
troppo generiche. In altri casi il modello lightweight aggiungeva meccanicamente termini
istituzionali come ``DIEM'', spostando il retrieval verso documenti non pertinenti. La
stabilizzazione ha portato al rewrite obbligatorio prima di ogni retrieve, al
retry-aware rewrite e alla risoluzione dei follow-up sulla history prima dello scope check.

\textbf{Tool e alternative non adottate.} Alcune soluzioni inizialmente promettenti sono
state scartate perché introducevano più fragilità che benefici: i tool ausiliari di
sintesi rendevano meno prevedibile il ciclo di retrieval, \code{calculate} aumentava la
latenza e dipendeva comunque da espressioni generate dal modello lightweight, mentre chunk
più grandi miglioravano l'autocontenimento ma aggiungevano rumore. Anche Crawl4AI è stato
testato come alternativa a trafilatura: poiché
le sorgenti DIEM sono esposte quasi sempre come HTML già renderizzato, il browser headless
estraeva poco contenuto realmente aggiuntivo e le prestazioni end-to-end restavano
comparabili, senza giustificare la maggiore complessità operativa.

Questa fase ha quindi orientato lo sviluppo verso interventi meno visibili ma più
determinanti: stabilità del ciclo rewrite$\to$retrieve, controllo della history,
prompting più vincolante e valutazione più robusta. Le configurazioni C1--C5 della
Sezione~\ref{sec:eval_results} misurano l'effetto cumulativo di queste decisioni.

%% ============================================================
\section{Valutazione Metriche e Risultati}
\label{sec:eval_results}

Il presente capitolo riporta i risultati della valutazione quantitativa condotta su
cinque configurazioni evolutive del sistema \sys{}. Le prime quattro (C1--C4)
introducono modifiche incrementali --- al modello linguistico, all'embedding, al reranker
o alla strategia di chunking --- per misurarne l'impatto sulle metriche Ragas e sui
criteri custom di scope awareness e robustezza. La quinta configurazione (C5)
rappresenta la versione finale del sistema: stessa infrastruttura di C3 con i fix
agentici successivi (rewrite retry-aware, sliding window,
regulation check, academic roles, altri miglioramenti di robustezza). La valutazione usa
il golden set italiano descritto nella Sezione~\ref{sec:evaluation}: 32~domande
\emph{in\_scope}/\emph{multi\_turn}, 8~domande \emph{out\_of\_scope} e 6~scenari di
robustezza. Il modello giudice è \code{meta-llama/llama-3.1-8b-instruct} per tutti i run.

\subsection{Configurazioni a Confronto}
\label{subsec:eval_configs}

Le cinque configurazioni rappresentano un percorso evolutivo da un sistema
completamente locale (C1) a una configurazione ibrida con componenti cloud per LLM,
embedding e reranking, con raffinamenti agentici incrementali (C5).
La Tabella~\ref{tab:configs} ne sintetizza i parametri principali.

\begin{table}[H]
\centering
\caption{Parametri delle cinque configurazioni valutate. \textit{BiK}: documenti
         estratti dal bi-encoder al primo stadio. \textit{CrossK}: documenti
         selezionati dopo il reranking. \textit{Parent\,/\,Child}: dimensione e overlap
         dei chunk in caratteri.}
\label{tab:configs}
\footnotesize
\setlength{\tabcolsep}{3.5pt}
\begin{tabular}{@{}cp{2.7cm}p{3.0cm}p{2.7cm}ccc@{}}
\toprule
 & \textbf{LLM} & \textbf{Embedding} & \textbf{Reranker}
 & \textbf{BiK / CrossK} & \textbf{Soglia} & \textbf{Parent\,/\,Child (car.)} \\
\midrule
\textbf{C1}
  & Qwen2.5-7B (locale)
  & e5-base 768d (locale)
  & cross-encoder locale
  & 20\,/\,3 & 0.50 & 2000\,/\,200\,\textbar\;400\,/\,50 \\
\addlinespace[3pt]
\textbf{C2}
  & GPT-oss-120B (cloud)
  & e5-base 768d (locale)
  & Cohere Rerank-4-Pro
  & 15\,/\,5 & 0.50 & 2000\,/\,200\,\textbar\;400\,/\,50 \\
\addlinespace[3pt]
\textbf{C3}
  & GPT-oss-120B (cloud)
  & Qwen3-Emb-8B 4096d (cloud)
  & Cohere Rerank-4-Pro
  & 15\,/\,5 & 0.50 & 2000\,/\,200\,\textbar\;400\,/\,50 \\
\addlinespace[3pt]
\textbf{C4}
  & GPT-oss-120B (cloud)
  & Qwen3-Emb-8B 4096d (cloud)
  & Cohere Rerank-4-Pro
  & 15\,/\,5 & 0.50 & 3750\,/\,375\,\textbar\;600\,/\,80 \\
\addlinespace[3pt]
\textbf{C5}
  & GPT-oss-120B (cloud)
  & Qwen3-Emb-8B 4096d (cloud)
  & Cohere Rerank-4-Pro
  & 15\,/\,5 & 0.50 & 2000\,/\,200\,\textbar\;400\,/\,50 \\
\bottomrule
\end{tabular}
\end{table}

\noindent La soglia \code{RETRIEVER\_SCORE\_THRESHOLD} è rimasta invariata a 0,50 in
tutte le configurazioni.

Le transizioni sono quattro. C1$\to$C2: LLM locale sostituito da un modello cloud di
larga scala, con reranker Cohere Rerank-4-Pro e CrossK aumentato
a 5. C2$\to$C3: embedding locale a 768d sostituito da un modello cloud ad alta
dimensionalità (4096d), con re-indicizzazione completa su un nuovo Chroma store.
C3$\to$C4: chunk padre $+87.5\%$ (2000$\to$3750 caratteri) e chunk figlio $+50\%$
(400$\to$600), per contesti più autocontenuti. C4$\to$C5: ritorno al chunking di C3
(2000/200 | 400/50) con i raffinamenti agentici descritti nella
Sezione~\ref{sec:core-rag} (rewrite retry-aware, sliding window, verifica del regolamento
di laurea, precedenza accademica), che agiscono sul comportamento del grafo senza
alterare il retrieval.

\subsection{Risultati Aggregati}
\label{subsec:eval_results_table}

La Tabella~\ref{tab:allresults} mostra i valori di tutte le metriche per le cinque
configurazioni. I valori in \textbf{grassetto} indicano il massimo per ciascuna riga.
Nessuna configurazione domina in modo assoluto: ogni transizione migliora alcune
metriche a scapito di altre, evidenziando i compromessi tipici dei sistemi RAG.

\begin{table}[H]
\centering
\caption{Valori aggregati per configurazione. \textit{Scope Awareness} e
         \textit{Robustness} sono pass rate rispettivamente su 8 e 6 scenari;
         le restanti metriche sono score Ragas $\in[0,1]$. In parentesi il numero
         di domande con valutazione valida. $\dagger$ valori corretti con re-run
         selettivo (si veda il testo).}
\label{tab:allresults}
\small
\begin{tabular}{@{}lccccc@{}}
\toprule
\textbf{Metrica} & \textbf{C1} & \textbf{C2} & \textbf{C3} & \textbf{C4} & \textbf{C5} \\
\midrule
\rowcolor{blue!4}
ResponseRelevancy   & 0.703 \scriptsize(32/32) & \textbf{0.877} \scriptsize(32/32)
                    & 0.681 \scriptsize(32/32) & 0.717 \scriptsize(32/32)
                    & 0.682 \scriptsize(32/32) \\
Faithfulness        & \textbf{0.764} \scriptsize(32/32) & 0.664 \scriptsize(31/32)
                    & 0.696 \scriptsize(30/32) & 0.675 \scriptsize(31/32)
                    & 0.752$^\dagger$ \scriptsize(32/32) \\
\rowcolor{blue!4}
AnswerCorrectness   & 0.537 \scriptsize(31/32) & \textbf{0.603} \scriptsize(30/32)
                    & 0.527 \scriptsize(31/32) & 0.597 \scriptsize(30/32)
                    & 0.582 \scriptsize(32/32) \\
LLMContextPrecision & \textbf{0.711} \scriptsize(32/32) & 0.526 \scriptsize(32/32)
                    & 0.584 \scriptsize(32/32) & 0.667 \scriptsize(31/32)
                    & 0.633 \scriptsize(32/32) \\
\rowcolor{blue!4}
LLMContextRecall    & 0.737 \scriptsize(30/32) & 0.797 \scriptsize(32/32)
                    & 0.820 \scriptsize(32/32) & 0.799 \scriptsize(29/32)
                    & \textbf{0.828} \scriptsize(32/32) \\
Coherence           & 0.737 \scriptsize(19/32) & 0.815 \scriptsize(27/32)
                    & 0.815 \scriptsize(27/32) & 0.750 \scriptsize(20/32)
                    & \textbf{0.826}$^\dagger$ \scriptsize(32/32) \\
\rowcolor{blue!4}
Scope Awareness     & 0.375 \scriptsize(3/8)  & 0.500 \scriptsize(4/8)
                    & \textbf{0.750} \scriptsize(6/8)  & 0.500 \scriptsize(4/8)
                    & 0.625 \scriptsize(5/8) \\
Robustness          & 0.333 \scriptsize(2/6)  & \textbf{0.833} \scriptsize(5/6)
                    & 0.667 \scriptsize(4/6)  & 0.667 \scriptsize(4/6)
                    & 0.667 \scriptsize(4/6) \\
\bottomrule
\end{tabular}
\end{table}

\subsection{Analisi per Metrica}
\label{subsec:per_metric}

Sulle metriche generative, il salto da C1 a C2 (LLM cloud) porta guadagni netti su
ResponseRelevancy (0.703$\to$0.877, $+24.8\%$) e AnswerCorrectness (0.537$\to$0.603,
$+12.3\%$), grazie a risposte più focalizzate e accurate; il calo in C3 su entrambe
(0.681 e 0.527) è attribuibile alla re-indicizzazione su un nuovo Chroma store con
embedding Qwen3, che nella fase iniziale produce contesti meno centrati rispetto a un
corpus ottimizzato iterativamente, mentre C4 recupera parzialmente (0.717 e 0.597)
grazie a chunk più autocontenuti. Sulla Faithfulness, C1 (LLM locale, 0.764) resta la
configurazione più fedele al contesto: i modelli cloud (C2--C4, 0.664--0.696) tendono a
integrare conoscenza pretrainata anche in presenza di contesto esplicito.

Sulla qualità del contesto, LLMContextPrecision e LLMContextRecall mostrano traiettorie
complementari: la soglia più selettiva di C1 massimizza la precisione (0.711) ma
penalizza la recall (0.737), mentre l'abbassamento della soglia e l'aumento di CrossK in
C2--C4, insieme all'embedding Qwen3 ad alta dimensionalità, spostano il compromesso verso
la recall, con C3 al picco (0.820) e minimo di precisione in C2 (0.526). La Coherence
segue un comportamento a due livelli, 0.737 con LLM locale (C1) contro 0.815 con LLM
cloud (C2/C3, $+10.6\%$), con lieve regressione in C4 (0.750) attribuibile a contesti più
lunghi che introducono ridondanza nella sintesi; il campione valido per questa metrica è
però ridotto (19--27 domande su 32), il che ne limita la significatività statistica.

Sui criteri custom, Scope Awareness e Robustness seguono traiettorie non monotone. Scope
Awareness passa da 37.5\% (C1) a un picco di 75\% in C3 --- l'embedding Qwen3 distingue
meglio contenuto DIEM da conoscenza generica --- per poi regredire a 50\% in C4, dove
chunk più grandi aumentano la probabilità di contenuto parzialmente fuori dominio.
Robustness compie il salto più marcato dell'intera sperimentazione, da 33.3\% (C1) a
83.3\% (C2, $+150\%$), per poi stabilizzarsi al 66.7\% da C3 in poi: variazioni di
chunking e raffinamenti agentici non influiscono più in modo significativo su questa
metrica.

\subsubsection{Configurazione finale (C5)}

C5 condivide la stessa infrastruttura di retrieval di C3 (embedding Qwen3-Emb-8B,
chunking 2000/200 | 400/50) e introduce i miglioramenti agentici descritti nel
core RAG della Sezione~\ref{sec:core-rag}. Dal punto di vista delle metriche, i risultati
più rilevanti sono:
\begin{itemize}
  \item \textbf{LLMContextRecall 0.828}: nuovo massimo assoluto ($+1.0\%$ su C3). Il
    principale collo di bottiglia era il nodo \emph{rewrite}: analizzando i retrieval
    con \code{scripts/debug\_retrieval.py} è emerso che le riformulazioni del modello
    lightweight aggiungevano meccanicamente termini generici (``DIEM'') che spostavano
    il risultato su documenti non pertinenti. La revisione del prompt di rewrite, il
    meccanismo retry-aware e la forzatura dell'ordine rewrite$\to$retrieve hanno reso il
    comportamento controllato e prevedibile nella maggior parte dei casi.
  \item \textbf{Coherence 0.826$^\dagger$}: nuovo massimo assoluto ($+1.4\%$ su C2/C3),
    per due ragioni: tecnica (il re-run porta il campione valido da 27/32 a 32/32) e
    qualitativa (il retrieval più stabile di C5 produce contesti più pertinenti anche
    su domande singole).
  \item \textbf{Faithfulness 0.752$^\dagger$}: valore più elevato tra le configurazioni
    cloud (C1 locale resta prima con 0.764), grazie a regole di citazione esplicite nel
    prompt e alla rimozione del tool \code{calculate}.
  \item \textbf{Scope Awareness 0.625 (5/8)}: apparentemente inferiore a C3 (0.750), ma
    tutti e tre i fallimenti sono rifiuti corretti valutati erroneamente come scorretti
    dal giudice automatico.
\end{itemize}

\textbf{Nota operativa sulla demo.} Un vector store aggiornato all'ultima settimana di
sviluppo ha mostrato un lieve calo di performance rispetto allo store consolidato
(instabilità dei context header tra run, prompt/soglie ottimizzati sullo store
precedente): per la demo finale si è quindi privilegiata riproducibilità usando lo store
consolidato.

\subsection{Affidabilità del Giudice Automatico e Re-run Correttivo}
\label{subsec:ragas_rerun}

Ragas usa un LLM come giudice per \emph{Coherence} e \emph{Faithfulness}, con due livelli
di incertezza: \emph{strutturale} (il giudice può sbagliare la valutazione --- il caso
della Scope Awareness di C5, §\ref{subsec:per_metric}, dove tutti i fallimenti sono
rifiuti corretti valutati erroneamente) e \emph{tecnico} (su risposte lunghe il giudice
può terminare anticipatamente o produrre output non parsabili, restituendo \texttt{NaN} e
abbassando artificialmente la media). Per C5, il run iniziale ha prodotto valori
depressi (Coherence 0.526, Faithfulness 0.578) per questo secondo motivo; lo script
\code{evaluation/rerun\_ragas.py} ha corretto il problema reinvocando il giudice solo
sulle coppie \texttt{NaN} (senza rigenerare risposte o ri-eseguire la pipeline), portando
i valori a Coherence 0.826 e Faithfulness 0.752 ($\dagger$ in Tabella~\ref{tab:allresults}).
Il re-run correttivo è stato applicato solo a queste due metriche perché lì si sono
manifestati errori tecnici di parsing/terminazione; le altre metriche sono state mantenute
come prodotte dal runner originale. Poiché diverse metriche Ragas dipendono comunque da
valutazione automatica, i risultati vanno letti come indicatori comparativi e non come
ground truth assoluta: la scelta di C5 come configurazione finale è stata quindi
supportata anche da test manuali sistematici.

\subsection{Riepilogo dell'Evoluzione}
\label{subsec:evolution_summary}

Le Figure~\ref{fig:ragas_metrics} e~\ref{fig:scope_robust} visualizzano l'evoluzione
delle metriche Ragas e dei criteri custom sulle cinque configurazioni.

\begin{figure}[H]
\centering
\begin{tikzpicture}
\begin{axis}[
  ybar,
  bar width=7pt,
  width=0.96\textwidth,
  height=5.5cm,
  enlarge x limits=0.10,
  legend style={
    at={(0.5,1.03)}, anchor=south,
    legend columns=5,
    font=\scriptsize,
    /tikz/every even column/.append style={column sep=0.25em}
  },
  ylabel={Punteggio},
  ymin=0, ymax=1.05,
  ytick={0,0.2,0.4,0.6,0.8,1.0},
  xtick={1,2,3,4,5,6},
  xticklabels={{Resp.\ Relevancy},{Faithfulness},{Ans.\ Correctness},
               {Ctx.\ Precision},{Ctx.\ Recall},{Coherence}},
  xticklabel style={font=\footnotesize, rotate=28, anchor=east},
  grid=major,
  grid style={dotted, gray!35},
]
\addplot[fill=blue!60, draw=blue!80!black] coordinates
  {(1,0.703)(2,0.764)(3,0.537)(4,0.711)(5,0.737)(6,0.737)};
\addplot[fill=orange!65, draw=orange!80!black] coordinates
  {(1,0.877)(2,0.664)(3,0.603)(4,0.526)(5,0.797)(6,0.815)};
\addplot[fill=green!50!black!70, draw=green!60!black] coordinates
  {(1,0.681)(2,0.696)(3,0.527)(4,0.584)(5,0.820)(6,0.815)};
\addplot[fill=purple!55, draw=purple!75!black] coordinates
  {(1,0.717)(2,0.675)(3,0.597)(4,0.667)(5,0.799)(6,0.750)};
\addplot[fill=red!55!black!70, draw=red!75!black] coordinates
  {(1,0.682)(2,0.752)(3,0.582)(4,0.633)(5,0.828)(6,0.826)};
\legend{C1 (baseline locale), C2 (LLM cloud), C3 (embed.\ cloud), C4 (chunk ottim.), C5 (fix agentici)}
\end{axis}
\end{tikzpicture}
\caption{Metriche Ragas per le cinque configurazioni. Ogni gruppo di cinque barre
         corrisponde a una metrica; i colori distinguono le configurazioni.
         I valori di Faithfulness e Coherence per C5 sono corretti post re-run ($\dagger$).}
\label{fig:ragas_metrics}
\end{figure}

\begin{figure}[H]
\centering
\begin{tikzpicture}
\begin{axis}[
  ybar,
  bar width=12pt,
  width=0.65\textwidth,
  height=4.8cm,
  enlarge x limits=0.40,
  legend style={
    at={(1.05,0.5)}, anchor=west,
    legend columns=1,
    font=\small
  },
  ylabel={Pass rate},
  ymin=0, ymax=1.05,
  ytick={0,0.2,0.4,0.6,0.8,1.0},
  xtick={1,2},
  xticklabels={{Scope Awareness},{Robustness}},
  xticklabel style={font=\small},
  grid=major,
  grid style={dotted, gray!35},
]
\addplot[fill=blue!60, draw=blue!80!black] coordinates {(1,0.375)(2,0.333)};
\addplot[fill=orange!65, draw=orange!80!black] coordinates {(1,0.500)(2,0.833)};
\addplot[fill=green!50!black!70, draw=green!60!black] coordinates {(1,0.750)(2,0.667)};
\addplot[fill=purple!55, draw=purple!75!black] coordinates {(1,0.500)(2,0.667)};
\addplot[fill=red!55!black!70, draw=red!75!black] coordinates {(1,0.625)(2,0.667)};
\legend{C1, C2, C3, C4, C5}
\end{axis}
\end{tikzpicture}
\caption{Pass rate per Scope Awareness (8~scenari out-of-scope) e Robustness
         (6~scenari avversariali) nelle cinque configurazioni.}
\label{fig:scope_robust}
\end{figure}

Il passaggio dal baseline locale (C1) alla configurazione finale (C5) produce
guadagni netti sostanziali su Robustness ($+100\%$, da 33.3\% a 66.7\%),
LLMContextRecall ($+12.3\%$, da 0.737 a 0.828, nuovo massimo assoluto) e
Coherence ($+12.1\%$, da 0.737 a 0.826$^\dagger$). A fronte di questi guadagni
si osserva una regressione su Faithfulness ($-1.6\%$, da 0.764 a 0.752$^\dagger$)
e LLMContextPrecision ($-11.0\%$, da 0.711 a 0.633), entrambe spiegabili con
l'allargamento del retrieval (soglia più bassa, CrossK maggiore) e la tendenza
dei modelli di larga scala a integrare conoscenza esterna --- un compromesso
accettabile in un sistema orientato alla copertura e alla robustezza conversazionale.
C5, come configurazione finale, raggiunge il miglior equilibrio globale tra copertura,
coerenza e affidabilità agentistica.

%% ============================================================
\section{Valutazione Metriche e Risultati}
\label{sec:eval_results}

Il presente capitolo riporta i risultati della valutazione quantitativa condotta su
quattro configurazioni evolutive del sistema \sys{}. Ciascuna configurazione introduce
modifiche incrementali --- al modello linguistico, all'embedding, al reranker o alla
strategia di chunking --- per misurarne l'impatto sulle metriche Ragas e sui criteri
custom di scope awareness e robustezza. Il golden set utilizzato è
\code{golden\_set\_it.json}: 32~domande nelle categorie \emph{in\_scope} e
\emph{multi\_turn}, 8~domande \emph{out\_of\_scope} e 6~scenari di robustezza.
Il modello giudice è \code{meta-llama/llama-3.1-8b-instruct} per tutti e quattro i run.

\subsection{Configurazioni a Confronto}
\label{subsec:eval_configs}

Le quattro configurazioni rappresentano un percorso evolutivo da un sistema
completamente locale (C1) a una configurazione ibrida con componenti cloud per LLM,
embedding e reranking (C4). La Tabella~\ref{tab:configs} ne sintetizza i parametri
principali.

\begin{table}[H]
\centering
\caption{Parametri delle quattro configurazioni valutate. \textit{BiK}: documenti
         estratti dal bi-encoder al primo stadio. \textit{CrossK}: documenti
         selezionati dopo il reranking. \textit{Parent\,/\,Child}: dimensione e overlap
         dei chunk in token.}
\label{tab:configs}
\footnotesize
\setlength{\tabcolsep}{3.5pt}
\begin{tabular}{@{}cp{2.7cm}p{3.0cm}p{2.7cm}ccc@{}}
\toprule
 & \textbf{LLM} & \textbf{Embedding} & \textbf{Reranker}
 & \textbf{BiK / CrossK} & \textbf{Soglia} & \textbf{Parent\,/\,Child (token)} \\
\midrule
\textbf{C1}
  & Qwen2.5-7B (locale)
  & e5-base 768d (locale)
  & cross-encoder locale
  & 20\,/\,3 & 0.50 & 2000\,/\,200\,\textbar\;400\,/\,50 \\
\addlinespace[3pt]
\textbf{C2}
  & GPT-oss-120B (cloud)
  & e5-base 768d (locale)
  & Cohere Rerank-4-Pro
  & 15\,/\,5 & 0.45 & 2000\,/\,200\,\textbar\;400\,/\,50 \\
\addlinespace[3pt]
\textbf{C3}
  & GPT-oss-120B (cloud)
  & Qwen3-Emb-8B 4096d (cloud)
  & Cohere Rerank-4-Pro
  & 15\,/\,5 & 0.45 & 2000\,/\,200\,\textbar\;400\,/\,50 \\
\addlinespace[3pt]
\textbf{C4}
  & GPT-oss-120B (cloud)
  & Qwen3-Emb-8B 4096d (cloud)
  & Cohere Rerank-4-Pro
  & 15\,/\,5 & 0.45 & 3750\,/\,375\,\textbar\;600\,/\,80 \\
\bottomrule
\end{tabular}
\end{table}

Le variazioni introdotte ad ogni transizione sono tre. C1$\to$C2: sostituzione del
LLM locale con un modello cloud di larga scala e adozione del reranker Cohere
Rerank-4-Pro, con soglia di recupero abbassata da 0.50 a 0.45 e CrossK aumentato
da 3 a 5. C2$\to$C3: sostituzione dell'embedding locale a 768 dimensioni con un
modello cloud ad alta dimensionalità (4096d) e re-indicizzazione completa del corpus
su un nuovo Chroma store. C3$\to$C4: aumento dei chunk padre del $+87.5\%$
(da 2000 a 3750 token) e dei chunk figlio del $+50\%$ (da 400 a 600 token), con
overlap scalato proporzionalmente, per fornire al modello contesti più autocontenuti.

\subsection{Risultati Aggregati}
\label{subsec:eval_results_table}

La Tabella~\ref{tab:allresults} mostra i valori di tutte le metriche per le quattro
configurazioni. I valori in \textbf{grassetto} indicano il massimo per ciascuna riga.
Nessuna configurazione domina in modo assoluto: ogni transizione migliora alcune
metriche a scapito di altre, evidenziando i compromessi tipici dei sistemi RAG.

\begin{table}[H]
\centering
\caption{Valori aggregati per configurazione. \textit{Scope Awareness} e
         \textit{Robustness} sono pass rate rispettivamente su 8 e 6 scenari;
         le restanti metriche sono score Ragas $\in[0,1]$. In parentesi il numero
         di domande con valutazione valida.}
\label{tab:allresults}
\small
\begin{tabular}{@{}lcccc@{}}
\toprule
\textbf{Metrica} & \textbf{C1} & \textbf{C2} & \textbf{C3} & \textbf{C4} \\
\midrule
\rowcolor{blue!4}
ResponseRelevancy   & 0.703 \scriptsize(32/32) & \textbf{0.877} \scriptsize(32/32)
                    & 0.681 \scriptsize(32/32) & 0.717 \scriptsize(32/32) \\
Faithfulness        & \textbf{0.764} \scriptsize(32/32) & 0.664 \scriptsize(31/32)
                    & 0.696 \scriptsize(30/32) & 0.675 \scriptsize(31/32) \\
\rowcolor{blue!4}
AnswerCorrectness   & 0.537 \scriptsize(31/32) & \textbf{0.603} \scriptsize(30/32)
                    & 0.527 \scriptsize(31/32) & 0.597 \scriptsize(30/32) \\
LLMContextPrecision & \textbf{0.711} \scriptsize(32/32) & 0.526 \scriptsize(32/32)
                    & 0.584 \scriptsize(32/32) & 0.667 \scriptsize(31/32) \\
\rowcolor{blue!4}
LLMContextRecall    & 0.737 \scriptsize(30/32) & 0.797 \scriptsize(32/32)
                    & \textbf{0.820} \scriptsize(32/32) & 0.799 \scriptsize(29/32) \\
Coherence           & 0.737 \scriptsize(19/32) & \textbf{0.815} \scriptsize(27/32)
                    & \textbf{0.815} \scriptsize(27/32) & 0.750 \scriptsize(20/32) \\
\rowcolor{blue!4}
Scope Awareness     & 0.375 \scriptsize(3/8)  & 0.500 \scriptsize(4/8)
                    & \textbf{0.750} \scriptsize(6/8)  & 0.500 \scriptsize(4/8) \\
Robustness          & 0.333 \scriptsize(2/6)  & \textbf{0.833} \scriptsize(5/6)
                    & 0.667 \scriptsize(4/6)  & 0.667 \scriptsize(4/6) \\
\bottomrule
\end{tabular}
\end{table}

\subsection{Analisi per Metrica}
\label{subsec:per_metric}

\subsubsection{Rilevanza della risposta (ResponseRelevancy)}

Il salto da C1 (0.703) a C2 (0.877, $+24.8\%$) evidenzia come il modello cloud
sia nettamente più efficace nel produrre risposte focalizzate sulla domanda, senza
divagazioni o risposte parziali. Il calo da C2 a C3 (0.681, $-22.3\%$) è attribuibile
alla re-indicizzazione su un nuovo Chroma store con embedding Qwen3: nella fase
iniziale il nuovo indice può produrre contesti meno centrati rispetto a un corpus
ottimizzato iterativamente. C4 recupera a 0.717 ($+5.3\%$ su C3), beneficiando
di chunk più autocontenuti che limitano il rumore nel contesto fornito al generatore.

\subsubsection{Fedeltà al contesto (Faithfulness)}

C1 registra la Faithfulness più elevata (0.764): il modello locale, avendo capacità
generative più limitate, tende a restare più ancorato al testo estratto. Le
configurazioni cloud (C2--C4, valori 0.664--0.696) mostrano valori inferiori,
coerentemente con la propensione dei modelli di larga scala a integrare conoscenza
pretrainata anche in presenza di un contesto esplicito. La lieve risalita da C2 a
C3 ($+4.8\%$) è coerente con il miglioramento qualitativo del retrieval: contesti
più rilevanti riducono la necessità del modello di ricorrere a fonti esterne.

\subsubsection{Correttezza della risposta (AnswerCorrectness)}

L'AnswerCorrectness segue una traiettoria analoga alla ResponseRelevancy: picco in
C2 (0.603, $+12.3\%$ su C1), flessione in C3 (0.527) per le stesse ragioni legate
alla re-indicizzazione, e quasi completo recupero in C4 (0.597). Il modello cloud
è strutturalmente più accurato nell'estrapolare fatti dal contesto e nel formulare
risposte complete rispetto al riferimento atteso nel golden set.

\subsubsection{Qualità del contesto recuperato}

\textbf{LLMContextPrecision.} Il baseline C1 mantiene la precisione più alta (0.711),
grazie alla soglia più selettiva (0.50) e al numero ridotto di chunk conservati
(CrossK=3). L'abbassamento della soglia e l'aumento di CrossK in C2--C4 incrementano
la copertura a scapito della precisione (minimo in C2: 0.526, $-26.0\%$).
Il trend di recupero C2$\to$C3$\to$C4 (0.526$\to$0.584$\to$0.667) dimostra che
l'embedding Qwen3 ad alta dimensionalità produce corrispondenze semantiche più
accurate e che chunk più lunghi, pur in numero minore, contengono informazioni più
pertinenti nel loro insieme.

\textbf{LLMContextRecall.} La recall mostra la traiettoria complementare e opposta:
miglioramento progressivo da C1 (0.737) al picco in C3 (0.820, $+11.3\%$ su C1),
con lieve calo in C4 (0.799). L'embedding ad alta dimensionalità (4096d) è più
efficace nel catturare la semantica completa del corpus italiano multidominio.
Il calo in C4 è spiegabile con la riduzione del numero totale di chunk indicizzati
derivante da frammenti più grandi, con conseguente minore ridondanza nel corpus.

\subsubsection{Coerenza della risposta (Coherence)}

La Coherence presenta un netto comportamento a due livelli: 0.737 con il LLM locale
(C1) e 0.815 con il LLM cloud (C2 e C3, $+10.6\%$). L'intero guadagno è attribuibile
alla qualità generativa superiore del modello cloud, che produce risposte strutturalmente
più coese e logicamente organizzate. La lieve regressione in C4 (0.750) è compatibile
con l'ipotesi che contesti più lunghi, pur informativi, possano introdurre ridondanza
che rende la sintesi della risposta meno fluida. Si segnala che il campione di valutazione
per la Coherence è ridotto rispetto alle altre metriche (19--20 domande valide su 32
per C1 e C4; 27/32 per C2 e C3), il che riduce la significatività statistica del
confronto su questa metrica.

\subsubsection{Consapevolezza dell'ambito (Scope Awareness)}

Il tasso di rifiuto corretto delle domande fuori ambito segue una traiettoria non
monotona: C1 (37.5\%, 3/8) $\to$ C2 (50.0\%, 4/8) $\to$ C3 (75.0\%, 6/8) $\to$
C4 (50.0\%, 4/8). Il picco in C3 indica che l'embedding Qwen3, addestrato su un
corpus più ampio e con dimensionalità quattro volte superiore, distingue meglio
tra contenuto specifico del DIEM e conoscenza generica fuori ambito. La regressione
in C4 (50\%), nonostante l'uso dello stesso embedding, è probabilmente dovuta ai
chunk più grandi: frammenti più estesi aumentano la probabilità di includere
contenuto parzialmente fuori dominio, inducendo il modello a rispondere invece di
applicare il guardrail di rifiuto.

\subsubsection{Robustezza (Robustness)}

Il miglioramento più marcato dell'intera sperimentazione riguarda la robustezza:
C1 (33.3\%, 2/6) $\to$ C2 (83.3\%, 5/6, $+150\%$). Il modello cloud dimostra
superiorità netta nella gestione di scenari avversariali, domande ambigue e challenge
alle risposte precedenti. Il calo C2$\to$C3 (66.7\%, 4/6) coincide con il cambio
di store e di embedding: la differente distribuzione dei contesti recuperati può
alterare il comportamento del guardrail su certi pattern avversariali specifici.
C4 mantiene il valore di C3 (66.7\%), confermando che la variazione delle dimensioni
dei chunk non influisce in misura significativa sulla robustezza.

\subsection{Riepilogo dell'Evoluzione}
\label{subsec:evolution_summary}

Le Figure~\ref{fig:ragas_metrics} e~\ref{fig:scope_robust} visualizzano l'evoluzione
delle metriche Ragas e dei criteri custom sulle quattro configurazioni.

\begin{figure}[H]
\centering
\begin{tikzpicture}
\begin{axis}[
  ybar,
  bar width=9pt,
  width=0.96\textwidth,
  height=7.5cm,
  enlarge x limits=0.10,
  legend style={
    at={(0.5,-0.22)}, anchor=north,
    legend columns=4,
    font=\small,
    /tikz/every even column/.append style={column sep=0.4em}
  },
  ylabel={Punteggio},
  ymin=0, ymax=1.05,
  ytick={0,0.2,0.4,0.6,0.8,1.0},
  xtick={1,2,3,4,5,6},
  xticklabels={{Resp.\ Relevancy},{Faithfulness},{Ans.\ Correctness},
               {Ctx.\ Precision},{Ctx.\ Recall},{Coherence}},
  xticklabel style={font=\footnotesize, rotate=28, anchor=east},
  grid=major,
  grid style={dotted, gray!35},
]
\addplot[fill=blue!60, draw=blue!80!black] coordinates
  {(1,0.703)(2,0.764)(3,0.537)(4,0.711)(5,0.737)(6,0.737)};
\addplot[fill=orange!65, draw=orange!80!black] coordinates
  {(1,0.877)(2,0.664)(3,0.603)(4,0.526)(5,0.797)(6,0.815)};
\addplot[fill=green!50!black!70, draw=green!60!black] coordinates
  {(1,0.681)(2,0.696)(3,0.527)(4,0.584)(5,0.820)(6,0.815)};
\addplot[fill=purple!55, draw=purple!75!black] coordinates
  {(1,0.717)(2,0.675)(3,0.597)(4,0.667)(5,0.799)(6,0.750)};
\legend{C1 (baseline locale), C2 (LLM cloud), C3 (embed.\ cloud), C4 (chunk ottim.)}
\end{axis}
\end{tikzpicture}
\caption{Metriche Ragas per le quattro configurazioni. Ogni gruppo di quattro barre
         corrisponde a una metrica; i colori distinguono le configurazioni.}
\label{fig:ragas_metrics}
\end{figure}

\begin{figure}[H]
\centering
\begin{tikzpicture}
\begin{axis}[
  ybar,
  bar width=16pt,
  width=0.60\textwidth,
  height=6.0cm,
  enlarge x limits=0.40,
  legend style={
    at={(1.05,0.5)}, anchor=west,
    legend columns=1,
    font=\small
  },
  ylabel={Pass rate},
  ymin=0, ymax=1.05,
  ytick={0,0.2,0.4,0.6,0.8,1.0},
  xtick={1,2},
  xticklabels={{Scope Awareness},{Robustness}},
  xticklabel style={font=\small},
  grid=major,
  grid style={dotted, gray!35},
]
\addplot[fill=blue!60, draw=blue!80!black] coordinates {(1,0.375)(2,0.333)};
\addplot[fill=orange!65, draw=orange!80!black] coordinates {(1,0.500)(2,0.833)};
\addplot[fill=green!50!black!70, draw=green!60!black] coordinates {(1,0.750)(2,0.667)};
\addplot[fill=purple!55, draw=purple!75!black] coordinates {(1,0.500)(2,0.667)};
\legend{C1, C2, C3, C4}
\end{axis}
\end{tikzpicture}
\caption{Pass rate per Scope Awareness (8~scenari out-of-scope) e Robustness
         (6~scenari avversariali) nelle quattro configurazioni.}
\label{fig:scope_robust}
\end{figure}

Il passaggio dal baseline locale (C1) alla configurazione finale (C4) produce
guadagni netti sostanziali su Robustness ($+100\%$, da 33.3\% a 66.7\%),
Scope Awareness ($+33.3\%$, da 37.5\% a 50.0\%) e AnswerCorrectness ($+11.2\%$,
da 0.537 a 0.597), con LLMContextRecall in miglioramento dell'$8.4\%$ (da 0.737
a 0.799). A fronte di questi guadagni si osservano regressioni su Faithfulness
($-11.6\%$, da 0.764 a 0.675) e LLMContextPrecision ($-6.2\%$, da 0.711 a 0.667),
entrambe spiegabili con l'allargamento del retrieval (soglia più bassa, CrossK
maggiore) e la tendenza dei modelli di larga scala a integrare conoscenza esterna ---
un compromesso accettabile in un sistema orientato alla copertura e alla robustezza
conversazionale.

%% ============================================================
\section{Discussione}
\label{sec:discussion}

\subsection{Principali Risultati}

L'architettura agentica basata su LangGraph StateGraph si è dimostrata superiore alle
chain lineari LCEL su più dimensioni, motivando a posteriori la scelta progettuale.

\textbf{Retrieval adattivo con retry.} Una chain lineare esegue il retrieve una sola
volta indipendentemente dalla qualità del risultato; il grafo agentico consente invece
al modello di richiamare il retrieve più volte con query progressivamente diversificate,
fino a un budget controllato (\texttt{MAX\_RETRIEVE\_CALLS}), abilitando il meccanismo
retry-aware descritto nella Sezione~\ref{sec:core-rag}.

\textbf{Interazione sicura a più livelli.} Una chain standard non ha punti naturali per
i guardrail. Il grafo esplicita otto nodi, due dedicati esclusivamente alla sicurezza
(\texttt{input\_guard}, \texttt{output\_guard}), con \texttt{scope\_guard} interposto tra
input e retrieval per rifiutare domande fuori ambito con latenza minima, senza
coinvolgere il modello principale.

\textbf{Controllo e manutenibilità.} La separazione esplicita dei nodi permette di
intervenire chirurgicamente su singoli comportamenti (fix del \texttt{recursion\_limit},
del nodo \texttt{forced\_retrieve}, del prompt di \texttt{rewrite}, aggiunta della regola
di precedenza accademica) senza destabilizzare o alterare la pipeline di retrieval
sottostante.

Il sistema di context header ha prodotto un miglioramento qualitativo misurabile nella
precisione del retrieval su frammenti semanticamente ambigui, applicando l'header a ogni
chunk figlio senza duplicarlo nei parent restituiti al modello. La validazione intermedia
ha inoltre mostrato che i guadagni più rilevanti non venivano dal cambio di estrattore,
ma dalla stabilizzazione della pipeline RAG: retrieval più controllato, rewrite meno
generico, history più pulita e prompt più vincolanti.

La configurazione finale C5 ha confermato che il principale collo di bottiglia del
sistema, una volta consolidata l'infrastruttura di retrieval, era la stabilità del ciclo
rewrite$\to$retrieve: migliorarne la robustezza (Sezione~\ref{subsec:per_metric}) ha
prodotto il miglior LLMContextRecall e Coherence dell'intera sperimentazione.

\subsection{Limiti del Sistema}
\label{subsec:limiti_sistema}

\begin{itemize}
  \item \textbf{Assenza di metadati temporali.} I CMS istituzionali non espongono segnali
    di data leggibili automaticamente: il sistema non può distinguere un programma
    dell'a.a.\ 2024/2025 da uno del 2021/2022 se l'anno non è esplicito nell'URL/testo.
    Verificato empiricamente con \code{scripts/diagnose\_date\_tags.py} (30 pagine
    campionate, nessun segnale affidabile via \texttt{Last-Modified} o \texttt{<time>}).
    Limite strutturale del dominio, non superabile senza modifiche ai CMS.
  \item \textbf{Copertura della knowledge base.} Snapshot statico del sito: pagine
    generate solo via JavaScript senza sitemap entry e contenuti dietro login (area
    studenti, esse3) non sono indicizzabili.
  \item \textbf{Faithfulness.} Su domande senza risposta nella knowledge base, il
    modello tende a produrre risposte plausibili invece di \texttt{KNOWLEDGE\_GAP}. Il
    tag e la regola di output obbligatorio riducono il fenomeno ma non lo eliminano.
  \item \textbf{Robustezza multi-turno.} Nelle versioni iniziali il retrieval drift al
    secondo turno faceva sì che una challenge (``Sei sicuro?'') producesse un re-retrieve
    su query diversa invece di confermare la risposta precedente. Mitigato con tre
    interventi: \code{scope\_guard} riscrive i follow-up prima della classificazione,
    \code{agent} riconosce challenge brevi via pattern lessicale e forza
    \code{rewrite()}$\to$\code{retrieve()}, e \code{REWRITE\_PROMPT} trasforma la
    challenge nella topic query del turno precedente. Il riconoscimento resta
    intenzionalmente conservativo e può fallire su formulazioni lunghe o indirette.
  \item \textbf{Instabilità del modello di rewrite.} Il modello lightweight (8B) usato
    per riformulare la query è strutturalmente meno affidabile su query specializzate,
    polisemiche o molto brevi, con riformulazioni errate o generiche che degradano
    silenziosamente il retrieve. Retry e forzatura dell'output (C5) mitigano senza
    eliminare il fenomeno; un modello più capace per il solo nodo rewrite resta
    un'ottimizzazione aperta.
  \item \textbf{Auto-valutazione della pertinenza.} Lo Step~2 chiede all'agente di
    verificare se il contesto risponde alla domanda, ma la valutazione è eseguita dallo
    stesso LLM che userà il contesto: su casi \emph{topic-adjacent} (es.\ ``programma del
    corso di Ingegneria del Software'' che atterra su piano di studi invece che scheda
    insegnamento) il modello può essere troppo ottimistico. Servirebbero segnali
    strutturati più forti (context header più discriminanti) o un judge di pertinenza
    separato.
  \item \textbf{Completezza su enumerazioni lunghe.} I cinque chunk padre selezionati dal
    reranker possono non coprire liste molto estese (es.\ strumentazione del Laboratorio
    MIVIA), rendendo impossibile una risposta esaustiva in un turno. Mitigabile
    aumentando CrossK/dimensione chunk (a costo di rumore e latenza) o affinando la
    query verso una sottocategoria.
  \item \textbf{Latenza di risposta.} In uso interattivo la media è circa 30 secondi, con
    tempi maggiori per retrieval multiplo e ragionamento numerico (es.\ voto di laurea).
  \item \textbf{Limiti di lettura dei PDF.} Estrazione testuale, non OCR visuale: tabelle
    e PDF scansionati possono produrre testo disordinato o illeggibile, motivando OCR
    dedicato.
  \item \textbf{Query multi-topic.} Domande che combinano due richieste distinte (es.\
    elenco laboratori + attrezzature di uno specifico) portano \code{rewrite()} a
    privilegiare la parte più generale, con retrieve insufficiente per la seconda
    sotto-domanda; mitigabile suddividendo in due turni, o con un nodo di query
    decomposition non implementato nella versione finale.
\end{itemize}

\subsection{Lavori Futuri}

\begin{itemize}
  \item \textbf{Modelli più capaci per enrichment e rewrite.} I nodi di context header e
    rewrite usano un modello lightweight da 8B, fonte di instabilità su query ambigue
    (Sezione~\ref{subsec:limiti_sistema}); un modello più capace migliorerebbe
    direttamente la precisione del retrieval senza alterare l'infrastruttura.
  \item \textbf{Fine-tuning sul dominio DIEM.} Un fine-tuning supervisionato del modello
    agente su conversazioni reali, domande del golden set e correzioni manuali
    specializzerebbe lo stile di risposta e ridurrebbe gli errori su entità specifiche
    (docenti, corsi, regolamenti), senza sostituire il RAG.
  \item \textbf{Human feedback e apprendimento continuo.} Il sistema non apprende dalle
    interazioni: feedback espliciti (\emph{thumbs up/down}) e impliciti (follow-up di
    correzione) potrebbero alimentare un dataset di preferenze per RLHF/DPO, il
    riordino dei documenti recuperati o regression test mirati.
  \item \textbf{Hybrid search (BM25 + dense).} Una componente BM25 sparse affiancata al
    bi-encoder Qwen3 migliorerebbe il recall su query con termini tecnici esatti (codici
    corso, acronimi) non sempre catturati dalla sola similarità semantica.
  \item \textbf{OCR avanzato per PDF.} Modelli OCR vision-language dedicati (es.\
    GLM-OCR) consentirebbero di indicizzare PDF scansionati o a layout complesso oggi
    non accessibili all'estrattore testuale.
  \item \textbf{Interfaccia vocale.} L'integrazione di TTS (es.\ ElevenLabs API)
    amplierebbe l'accessibilità senza modifiche al core RAG.
  \item \textbf{Crawling con agent browser.} Qualora venissero incluse nuove sorgenti
    realmente dinamiche, un agent browser potrebbe integrare il crawler HTML statico per
    accedere a contenuto JavaScript, accordion e sezioni lazy-loaded non indicizzabili.
\end{itemize}

%% ============================================================
\section{Conclusioni}
\label{sec:conclusioni}

Questo lavoro ha presentato la progettazione, l'implementazione e la valutazione di
\sys{}, un sistema RAG agentico per il supporto informativo universitario applicato al
Dipartimento DIEM dell'Università di Salerno. Il percorso tecnico ha attraversato cinque
configurazioni evolutive --- da un baseline locale a un grafo LangGraph con otto nodi
specializzati e una pipeline di retrieval cloud --- motivate da esigenze concrete emerse
durante i test: rifiuto affidabile delle domande fuori ambito, gestione robusta della
conversazione multi-turno, stabilità del ciclo rewrite$\to$retrieve e fedeltà delle
risposte al contesto.

I contributi principali sono: (1) una pipeline di ingestion specifica per siti universitari
italiani con crawling resiliente, caching incrementale, filtro temporale multi-tier, parsing
strutturato mirato e context header semantici su ogni chunk figlio; (2) un'architettura
agentica con guardrail a cascata, dual-provider per embedding/reranker/LLM e budget
controllato di retrieval; (3) meccanismi di robustezza agentistica (rewrite retry-aware,
sliding window, regulation check, academic roles) che stabilizzano il comportamento
senza alterare l'infrastruttura di retrieval; (4) l'integrazione di dati real-time
(EasyCourse) nella knowledge base; (5) una validazione intermedia che ha trasformato
criticità osservate nelle prime run in decisioni architetturali mirate.

La configurazione finale (C5) raggiunge LLMContextRecall 0.828 e Coherence 0.826
--- entrambi i massimi assoluti dell'intera sperimentazione --- con Faithfulness 0.752
e Robustness 66.7\%. I limiti residui più rilevanti sono l'instabilità strutturale del
modello lightweight utilizzato nel nodo rewrite e l'impossibilità di garantire
completezza su contenuti enumerativi che superano la finestra di retrieval. Le
direzioni di sviluppo prioritarie identificate (modelli più capaci per rewrite,
hybrid search BM25+dense, OCR avanzato) indirizzano direttamente questi limiti.

%% ============================================================
\newpage
\phantomsection
\addcontentsline{toc}{section}{Riferimenti bibliografici}
\bibliographystyle{plainnat}
\begin{thebibliography}{99}

\bibitem[Anthropic, 2024]{anthropic2024contextual}
Anthropic (2024).
\newblock \emph{Contextual Retrieval}.
\newblock Technical blog post, Anthropic, September 2024.

\bibitem[Barbaresi, 2021]{barbaresi2021trafilatura}
Barbaresi, A. (2021).
\newblock Trafilatura: A web scraping library and command-line tool for text
  discovery and extraction.
\newblock In \emph{Proceedings of ACL 2021}, pages 122--129.

\bibitem[Bansal et al., 2019]{bansal2019systematic}
Bansal, H. and Khan, R. (2019).
\newblock A review paper on human computer interaction.
\newblock \emph{International Journal of Advanced Research in Computer Science
  and Software Engineering}, 8(4):53--56.

\bibitem[Es et al., 2023]{es2023ragas}
Es, S., James, J., Anke, L.~E., and Schockaert, S. (2023).
\newblock Ragas: Automated evaluation of retrieval augmented generation.
\newblock \emph{arXiv preprint arXiv:2309.15217}.

\bibitem[LangChain, 2023]{langchain2023parentdoc}
LangChain (2023).
\newblock \emph{Parent Document Retriever}.
\newblock LangChain documentation, version 0.3.

\bibitem[LangGraph, 2024]{langgraph2024}
LangGraph (2024).
\newblock \emph{LangGraph: Building Stateful, Multi-Actor Applications with LLMs}.
\newblock LangChain Inc., 2024.

\bibitem[Lewis et al., 2020]{lewis2020rag}
Lewis, P., Perez, E., Piktus, A., Petroni, F., Karpukhin, V., Goyal, N., K\"{u}ttler, H.,
  Lewis, M., Yih, W., Rockt\"{a}schel, T., Riedel, S., and Kiela, D. (2020).
\newblock Retrieval-augmented generation for knowledge-intensive NLP tasks.
\newblock In \emph{Advances in Neural Information Processing Systems},
  volume~33, pages 9459--9474.

\bibitem[Nogueira \& Cho, 2020]{nogueira2020passage}
Nogueira, R. and Cho, K. (2020).
\newblock Passage re-ranking with BERT.
\newblock \emph{arXiv preprint arXiv:1901.04085}.

\bibitem[Shi et al., 2023]{shi2023replug}
Shi, W., Min, S., Yasunaga, M., Seo, M., James, R., Lewis, M., Zettlemoyer, L., and
  Yih, W. (2023).
\newblock REPLUG: Retrieval-augmented language model pre-training.
\newblock \emph{arXiv preprint arXiv:2301.12652}.

\bibitem[Smutny \& Schreiberova, 2020]{smutny2020chatbots}
Smutny, P. and Schreiberova, P. (2020).
\newblock Chatbots for learning: A review of educational chatbots for the Facebook
  Messenger.
\newblock \emph{Computers \& Education}, 151:103862.

\bibitem[OpenRouter, 2026]{openrouter2026gptoss120b}
OpenRouter (2026).
\newblock OpenAI: gpt-oss-120b -- API pricing \& benchmarks.
\newblock \url{https://openrouter.ai/openai/gpt-oss-120b}.

\bibitem[Suzgun \& Kalai, 2024]{suzgun2024metaprompting}
Suzgun, M. and Kalai, A.~T. (2024).
\newblock Meta-prompting: Enhancing language models with task-agnostic scaffolding.
\newblock \emph{arXiv preprint arXiv:2401.12954}.

\bibitem[Zheng et al., 2023]{zheng2023judging}
Zheng, L., Chiang, W.-L., Sheng, Y., Zhuang, S., Wu, Z., Zhuang, Y., Lin, Z.,
  Li, Z., Li, D., Xing, E.~P., Zhang, H., Gonzalez, J.~E., and Stoica, I. (2023).
\newblock Judging LLM-as-a-judge with MT-bench and chatbot arena.
\newblock \emph{arXiv preprint arXiv:2306.05685}.

\bibitem[Zheng et al., 2024]{zheng2023stepback}
Zheng, H., Mishra, S., Chen, X., Cheng, H.-T., Chi, E.~H., Le, Q.~V., and Zhou, D. (2024).
\newblock Take a step back: Evoking reasoning via abstraction in large language models.
\newblock \emph{International Conference on Learning Representations (ICLR)}.

\end{thebibliography}

\end{document}
